\chapter{具身RL仿真生态：mjlab、Isaac Lab 与 UniLab}\label{ch:rlmc-ecosystem}

% ════════════════════════════════════════════════════════════════
%  新部「强化学习运动控制」(P8_rl_motion_control) 的生态/选型实战章。
%  主线：以「仿真在 RL 管线中承担的工程职责」为概念骨架——仿真三角色、
%  GPU 并行、Manager 架构——再在每个概念下横向铺 mjlab / Isaac Lab / UniLab
%  三框架对比实战。假定读者已掌握 RL 理论（前面理论部 P6_rl 已讲 MDP/PPO/
%  Actor-Critic），不重教基础，聚焦工程扩展与上手。
%  源稿 RL运控/Ch01_生态系统.md（实践偏重）已按本主线重构；
%  关键 API/版本/benchmark/论文归属均联网核对（见各 \rebuilt / \pz / \cite）。
%  UniLab 列已据服务器实装框架的源码 + 实跑核对填实（异构 CPU-sim/GPU-policy，
%  论文 arXiv:2605.30313；非 Manager-Based）。凡 UniLab 真实欠缺处如实标注、不编造。
% ════════════════════════════════════════════════════════════════

强化学习运动控制（RL motion control）这一部，要把前几部的强化学习\emph{理论}——马尔可夫决策过程、贝尔曼方程、策略梯度、演员-评论家——落到一台真实（或仿真）的腿足/机械臂机器人上。而落地的第一块基础设施，不是奖励函数，也不是网络结构，而是\textbf{仿真器}：策略需要海量试错数据，真机无法承受这个量级的代价，于是仿真承担了「数据工厂」的角色。但「用哪个仿真器」在 2026 年已不是一个简单问题——MuJoCo、Isaac Gym、Isaac Lab、MuJoCo Playground、Newton、mjlab 构成了一个快速演进、互相借鉴的生态。本章是这一部的\textbf{地图章}：它不教你写某一行 reward，而是帮你建立全局视角，知道生态里有哪些框架、每一代解决了什么工程痛点、你的研究课题该选哪个。

本章遵循全书「概念/算法领路、工程实践兑现」的次序，但因为这是一章\emph{生态选型}内容，其「算法骨架」不是某个 RL 算法，而是\textbf{仿真在 RL 管线中承担的三类工程职责}：物理世界、MDP 机器、调试仪器。我们先把这三类职责讲清楚（它们决定了「一个仿真框架到底要提供什么」），再在每一类职责下横向铺开 mjlab、Isaac Lab、UniLab 三个框架的对比实战。读者会反复看到一条主线：\textbf{框架层是连接 RL 算法与物理引擎的桥梁}，掌握了框架层的设计方法论（无论是 mjlab/Isaac Lab 的 Manager-Based 体系，还是 UniLab 的 \texttt{NpEnv} 异构范式），就同时获得了对上层（算法接入）与下层（物理引擎切换）的控制力。

% ════════════════════════════════════════════════════════════════
\section*{环境配置指南}
\label{sec:rlmc-eco-setup}
\addcontentsline{toc}{section}{环境配置指南}

本章是概念/选型章，正文的代码以\emph{演示框架差异}为主，多为可读的最小片段而非完整工程。但要亲手验证后文每一处对比，你需要至少装好一个框架。下表给出两个主线框架的系统要求与版本兼容关系（数据以 2026 年中为快照，\textbf{版本号会随时间变化，请以官方仓库为准}）。

\begin{center}
\small
\begin{tabularx}{\linewidth}{@{}lLLL@{}}
\toprule
要素 & mjlab & Isaac Lab & 备注 \\
\midrule
操作系统 & Linux / macOS（CPU 调试）/ Win(WSL2) & Linux（推荐）/ Win & Isaac Sim 对 Linux 支持最完整 \\
Python & $\ge 3.10$（$<3.14$） & $\ge 3.10$ & 经核对 mjlab PyPI 元数据\,\cite{mjlab2026} \\
GPU & NVIDIA + CUDA（训练）& NVIDIA RTX 系（渲染）& mjlab 可 CPU 跑通流程做调试 \\
物理后端 & MuJoCo Warp & PhysX（默认）/ Newton(Beta) & 见 \cref{sec:rlmc-eco-evolution} \\
RL 后端 & RSL-RL（默认且唯一）& RSL-RL / RL Games / SKRL / SB3 & 见 \cref{sec:rlmc-eco-tools} \\
安装方式 & \Verb[breaklines]|pip install mjlab| / \Verb[breaklines]|uv sync| & conda + pip + Isaac Sim 或 kit-less & 复杂度差异显著 \\
启动延迟 & 数秒 & 数十秒；\textbf{首跑 $\ge 5$ 分钟}（编译 shader/USD）& 影响迭代体验 \\
\bottomrule
\end{tabularx}
\end{center}

\begin{note}[版本锚定：为什么要写死版本，又为什么不能迷信它]
机器人 RL 生态演进极快：2025 到 2026 年间，mjlab 从首个 Beta 发布到 1.x 系列，Isaac Lab 从 2.x 演进到 3.0 Beta，Newton 从公告走到 1.0 GA。若教材追最新版，读者在不同时间阅读会面对不同 API。本章的策略是\emph{锚定}一组主线版本（mjlab 1.x、Isaac Lab 2.3.2、RSL-RL 即 \texttt{rsl-rl-lib}，Isaac 钉 3.0.1 / mjlab 装 5.2.0），把 3.0 等破坏性变更放进注释框单独说明。\pz{源稿把 mjlab 锚定在 1.2.0；经核对 PyPI，截至 2026-05 最新已是 1.4.0\,\cite{mjlab2026}——故请把任何具体小版本号当作「最低兼容参考」，实际安装用当时的稳定版。}
\end{note}

\paragraph{分系统最小安装步骤（以 mjlab 为「先跑通」入口）。}mjlab 的安装是单条命令，适合作为第一个上手的框架；Isaac Lab 因依赖 Isaac Sim，步骤更长，建议在确认任务需要视觉渲染或多 RL 后端后再装（理由见 \cref{sec:rlmc-eco-select}）。

\begin{codebox}[bash]{mjlab:Linux / macOS / WSL2 通用最小安装}
# 推荐用 uv（Astral 的 Python 包管理器）管理环境，亦可纯 pip
pip install uv                      # 若已有 uv 可跳过
uv venv --python 3.11 && source .venv/bin/activate
uv pip install mjlab               # 一条命令装好;v1.0 起无需 git pin MuJoCo-Warp
python -c "import mjlab, importlib.metadata as m; print('mjlab', m.version('mjlab'))"  # mjlab 无 __version__ 属性，用 importlib.metadata 取版本
\end{codebox}

\begin{codebox}[bash]{Isaac Lab:Linux（pip 工作流，2.3.0 主线，简化示意）}
conda create -n isaaclab python=3.10 && conda activate isaaclab
pip install torch --index-url https://download.pytorch.org/whl/cu121
pip install isaacsim[all]          # Isaac Sim 运行时（数 GB，耗时较长）
git clone https://github.com/isaac-sim/IsaacLab.git
cd IsaacLab && ./isaaclab.sh --install   # 安装 isaaclab 系列扩展
# 验证:跑一个空场景，确认 Isaac Sim 能加载（首次启动会编译 shader，较慢）
\end{codebox}

第三个框架 UniLab 的安装介于两者之间：用 \Verb[breaklines]|uv sync| 选物理后端，不依赖 Isaac Sim（故无 mjlab 那条「别装 Omniverse」的坑），但物理在 CPU、需要一块 GPU 跑策略学习。

\begin{codebox}[bash]{UniLab:异构 CPU-sim + GPU-policy（uv 工作流，简化示意）}
git clone https://github.com/unilabsim/UniLab && cd UniLab
uv sync --extra mujoco              # MuJoCoUni 后端（mujoco-uni）;或 --extra motrix
# 国内首跑从 HuggingFace 拉资产，需换镜像端点:
export HF_ENDPOINT=https://hf-mirror.com
# 物理跑 CPU（核越多越好），策略需 GPU（CUDA / MPS / ROCm / XPU 任一）
\end{codebox}

\begin{pitfall}[安装阶段三个最常见的「卡住」]
\textbf{一·CUDA/PyTorch 版本不匹配}：Isaac Sim 对 CUDA 与 PyTorch 版本有强约束，装错会在第一次 \verb|import| 时报 \Verb[breaklines]|undefined symbol| 或段错误。务必照官方矩阵选版本，别用 \Verb[breaklines]|pip install torch| 的默认 wheel。\textbf{二·把 mjlab 当 Isaac Lab 那样装}：mjlab \emph{不需要} Isaac Sim/Omniverse，若你照 Isaac Lab 教程去装 Omniverse 再装 mjlab，纯属浪费时间。\textbf{三·无 GPU 也想跑训练}：mjlab 可在 CPU 上跑通\emph{流程}（调试用），但真正的并行训练需要 NVIDIA GPU；Isaac Lab 的渲染更依赖 RTX。先确认硬件再开始。
\end{pitfall}

\begin{pitfall}[两个 Isaac Lab 新手必踩的坑：启动慢、不能裸 Python 导入]
\textbf{一·首次启动远比「数秒」慢，要有心理预期。}上面那张系统要求表把 Isaac Lab 启动延迟写成「数十秒」，那是\emph{二次以后}的乐观值。实测在 RTX 4090 上跑最小的 \Verb[breaklines]|Isaac-Cartpole-v0|，从敲下命令到打印第 1 个训练迭代要 \textbf{$5$ 分钟以上}（纯花在 shader 编译 + USD 资产加载上）。这不是卡死：首跑会把编译结果\emph{缓存}下来，二次启动通常一两分钟。所以第一次跑 Isaac Lab 请用 \verb|--headless| 关掉渲染窗口、并预留 $5$--$8$ 分钟，别在三分钟时以为挂了就 \verb|Ctrl-C|（一旦中断缓存没写完，下次又得从头编译）。mjlab 无此问题（数秒起）。
\textbf{二·想「快速查个 API」时不能用系统 \texttt{python}。}Isaac Lab 的相机/资产类挂在 Omniverse 的 \verb|pxr|/\allowbreak\verb|omni.usd| 运行时上，必须先由 \verb|AppLauncher| 拉起 Isaac Sim 才能导入。直接 \Verb[breaklines]|python -c "import isaaclab.envs"| 会立刻 \Verb[breaklines]|ModuleNotFoundError: No module named 'pxr'|（或 \verb|omni.usd|）。正确姿势是用仓库自带的解释器包装器 \Verb[breaklines]|./isaaclab.sh -p your_script.py|，并在脚本里先 \Verb[breaklines]|AppLauncher(...)| 再 \verb|import|。这是「环境隔离」坑：报错信息看着像没装好，其实是没经 Isaac Sim 入口。
\end{pitfall}

% ════════════════════════════════════════════════════════════════
\section*{Quick Start：五分钟在 mjlab 上看见机器人}
\label{sec:rlmc-eco-quickstart}
\addcontentsline{toc}{section}{Quick Start：五分钟跑通}

装好 mjlab 后，下面四条命令让你在几分钟内完整体验「选环境 $\to$ 看默认姿态 $\to$ 看随机动作 $\to$ 启动训练」这条最小闭环。命令格式经核对为 mjlab 的 tyro CLI + \Verb[breaklines]|uv run|\,\cite{mjlab2026}；具体的任务注册名以你所装版本的 \verb|--help| 输出为准。

\begin{codebox}[bash]{mjlab 最小闭环（可逐条复制）}
# 1. 列出已注册任务，确认安装与注册都正常（无 ImportError）
uv run list-envs        # mjlab 列任务用 list-envs;train --help 只看 CLI 旗标
# 2. zero agent:动作恒为 0，机器人应保持默认站姿——验证模型正确加载
uv run play Mjlab-Velocity-Flat-Unitree-G1 --agent zero --num-envs 4
# 3. random agent:连续随机动作，应见明显运动但不「炸飞」、无 NaN
uv run play Mjlab-Velocity-Flat-Unitree-G1 --agent random --num-envs 4
# 4. 小规模训练 10 迭代:steps/s 稳定、reward 非 NaN 即说明管线打通
#    注:mjlab 默认用 WandB 记录;未配 API key 会报 "No API key configured".
#    冒烟可加 WANDB_MODE=disabled 跳过登录（正式训练再 wandb login）.
WANDB_MODE=disabled uv run train Mjlab-Velocity-Flat-Unitree-G1 \
    --env.scene.num-envs 512 --agent.max-iterations 10
\end{codebox}

\begin{insight}[zero agent 与 random agent 是「免费的体检」]\label{ins:rlmc-eco-zerorandom}
还没训练就先跑 \Verb[breaklines]|--agent zero| 与 \Verb[breaklines]|--agent random|，是机器人 RL 里一条极廉价、极有效的习惯。\textbf{zero agent} 把动作钉死为零：若机器人此时不是稳稳站着，而是瘫软或劈叉，说明默认姿态（offset）配错了——这与动作空间的「默认偏置」设计直接相关（详见 \cref{ch:rlmc-obsact} 的动作变换链）。\textbf{random agent} 把动作设为均匀随机：若机器人立刻飞出场景或数值变 NaN，多半是 action scale 太大或缺裁剪。两步都不写一行训练代码，却能在 30 秒内把「环境装错/配错」与「算法没调好」两类问题区分开——这正是 \cref{sec:rlmc-eco-roles} 将讲的仿真「调试仪器」角色的最小体现。
\end{insight}

% ════════════════════════════════════════════════════════════════
\section*{前置自测}
\label{sec:rlmc-eco-selftest}
\addcontentsline{toc}{section}{前置自测}

\begin{practice}[前置自测：答不出 $\ge 3$ 题，建议先回「强化学习理论部」复习]
\begin{enumerate}
  \item 强化学习的基本闭环里，观测、动作、奖励、下一观测四个量的因果关系是什么？（答错回看 \cref{sec:rl-basic-mdp}）
  \item \verb|sim.step()| 与 \verb|env.step()| 有何区别？后者在物理推进之外还做了哪些事？（答错回看 \cref{sec:rlmc-eco-roles}）
  \item PPO 的 rollout $\to$ GAE $\to$ clip 更新大致流程是什么？（本章不重讲，答错回看 \cref{ch:rl-pg} 与 \cref{ch:rl-ac}）
  \item 为什么「渲染漂亮」不等于「适合大规模 RL 训练」？请从吞吐量与显存两个角度回答。（答错回看 \cref{sec:rlmc-eco-gpu}）
  \item GPU 并行仿真与 CPU 多进程仿真的本质区别仅仅是「GPU 更快」吗？（答错回看 \cref{sec:rlmc-eco-gpu}）
\end{enumerate}
\end{practice}

\textbf{本章目标}：学完本章，你应当能够——
\begin{enumerate}
  \item \textbf{配置}并验证至少一个 GPU 并行 RL 框架（mjlab 或 Isaac Lab），跑通\nameref{sec:rlmc-eco-quickstart}的最小闭环；
  \item \textbf{区分} \verb|sim.step()| 与 \verb|env.step()| 的职责，定位一个 bug 是出在物理层还是 MDP 层；
  \item \textbf{对比} mjlab、Isaac Lab、UniLab 三框架在物理后端、安装复杂度、可视化、RL 后端、内置任务上的工程差异，并用决策树为一个新课题\textbf{选定}框架；
  \item \textbf{实现}一个粗略但够用的显存估算，判断「我的 GPU 能否跑 4096 并行环境」；
  \item \textbf{调通} Manager-Based 配置的最小骨架（observation/reward/event 三个 term），并解释它相对单体架构的工程收益；
  \item \textbf{定位} Isaac Lab 2.x 与 3.0 的破坏性变更（四元数序、warp-native 数据、模块改名），在升级报错时找到正确的迁移方向。
\end{enumerate}

% ════════════════════════════════════════════════════════════════
\section*{知识导航与阅读时间}
\label{sec:rlmc-eco-nav}
\addcontentsline{toc}{section}{知识导航}

\textbf{前置桥接}：本章假定你已学过强化学习理论部——尤其是马尔可夫决策过程（\cref{sec:rl-basic-mdp}）给出的 $(\mathcal S,\mathcal A,P,R,\gamma)$ 五元组、策略梯度（\cref{ch:rl-pg}）与演员-评论家（\cref{ch:rl-ac}）的 PPO/GAE 机制。本章\emph{不重讲}这些，而是把它们当作既有词汇，专注「这些算法要在什么样的仿真基础设施上运行」。本章也是本部入口章 \cref{ch:rlmc-obsact}（观测与动作设计）的\emph{前置}：先建立框架全景，再深入单个 Manager。

\begin{forest}
navtreeLR
[具身RL仿真生态
  [仿真的三类职责 (\S\ref*{sec:rlmc-eco-roles})
    [物理世界 sim.step]
    [MDP 机器 env.step]
    [调试仪器 ground-truth]
  ]
  [GPU 并行 (\S\ref*{sec:rlmc-eco-gpu})
    [数据零搬运]
    [SIMT / SoA]
    [显存与最大环境数]
  ]
  [生态演进六阶段 (\S\ref*{sec:rlmc-eco-evolution})
    [MuJoCo {$\to$} Isaac Gym]
    [Isaac Lab {$\to$} MJX]
    [Newton {$\to$} mjlab]
  ]
  [Manager 架构三框架对比 (\S\ref*{sec:rlmc-eco-manager})
    [mjlab]
    [Isaac Lab]
    [UniLab (非 Manager-Based)]
  ]
  [框架选型 (\S\ref*{sec:rlmc-eco-select})
    [决策树]
    [迁移成本]
  ]
  [其他框架与工具 (\S\ref*{sec:rlmc-eco-tools})]
]
\end{forest}

\begin{center}
\small
\begin{tabular}{@{}lll@{}}
\toprule
阅读方式 & 时间 & 适合谁 \\
\midrule
精读（含练习与动手） & 3--4 小时 & 初次接触机器人 RL 仿真的读者 \\
速读（跳过练习，看表与陷阱） & 1--1.5 小时 & 有 Isaac Gym/Gym 经验、不熟 mjlab/Newton 的读者 \\
速查（只看对比表与版本表） & 30 分钟 & 已有双框架经验，只查最新变化 \\
\bottomrule
\end{tabular}
\end{center}

% ════════════════════════════════════════════════════════════════
\section{概念骨架：仿真在 RL 管线中的三类职责}
\label{sec:rlmc-eco-roles}

\subsection{模块功能与在管线中的位置}

要回答「一个仿真框架到底该提供什么」，先要看清仿真在训练闭环里同时扮演的\emph{三个}角色。很多初学者只把仿真当成「一个跑得快的环境」，于是在调试时分不清问题出在哪一层。这三个角色，正是后文衡量 mjlab / Isaac Lab / UniLab 的三把尺子。

\begin{insight}[仿真之于机器人 RL，如编译器与单元测试之于软件工程——但边界要记牢]\label{ins:rlmc-eco-analogy}
你最终要把控制器跑在真机上（如同软件最终要上生产环境），但你不会把每一次语法错误都部署到生产系统才发现；同样，你不该把每一次奖励设计错误都拿真机摔一遍才知道。这个类比的\textbf{边界}是：编译器给出\emph{确定性}的错误信息，仿真器给出的却是对物理世界的\emph{近似}——「编译通过」约等于「语法正确」，但「仿真成功」\emph{不等于}「真机成功」。仿真能把试错成本压低几个数量级，却消除不了 sim-to-real gap。理解这条边界，是后续所有章节（尤其域随机化与 sim-to-real 部署）的前提。
\end{insight}

\paragraph{角色一·物理世界（物理引擎，\texttt{sim.step}）。}最底层职责：给定关节力矩，算出下一时刻的关节位置、速度与接触力。这由\emph{物理引擎}完成，而物理引擎的选择直接决定仿真的真实度——尤其是接触行为。下表是三类主流引擎的对照（接触模型差异将在 \cref{sec:rlmc-eco-evolution} 展开）：

\begin{center}
\small
\begin{tabular}{@{}llll@{}}
\toprule
物理引擎 & 坐标体系 & 接触模型 & 主要使用框架 \\
\midrule
MuJoCo / MuJoCo Warp & 广义坐标 & 凸优化（CG/Newton 求解器） & mjlab、MuJoCo Playground \\
PhysX & 笛卡尔坐标 & TGS 迭代求解器 & Isaac Lab（默认） \\
Newton & 多求解器统一 & MuJoCo Warp / Kamino / XPBD 等 & Isaac Lab 3.0 Beta \\
\bottomrule
\end{tabular}
\end{center}

\paragraph{角色二·MDP 机器（框架层，\texttt{env.step}）。}物理引擎只管 \verb|sim.step()|；而 RL 算法要消费的是一个完整的 MDP，这由框架层的 \verb|env.step()| 封装。\verb|env.step()| 在物理推进之外还做一整套事：

\begin{codebox}[text]{env.step(action) 在 sim.step() 之上封装的 MDP 逻辑}
env.step(action):
    1. 处理 action（缩放 scale / 裁剪 clip / 偏置 offset）
    2. 调用 sim.step() 一次或多次（action repeat / decimation）
    3. 计算 observation（从物理状态提取策略需要的信息）
    4. 计算 reward（评估当前行为）
    5. 判断 termination（是否结束当前 episode）
    6. 对终止的环境执行 reset
    7. 返回 (obs, reward, done, info)
\end{codebox}

\begin{insight}[\texttt{sim.step()} 与 \texttt{env.step()} 的区别，是理解整个框架架构的起点]\label{ins:rlmc-eco-simenv}
物理引擎负责「世界如何运转」，框架负责「如何把物理世界包装成 RL 算法可消费的 MDP」。把两者混为一谈，是初学者最常见的认知错误，后果是调试时分不清问题出在物理层还是 MDP 层：足端穿地、机器人「沉」进地下，多半是\emph{物理层}的接触参数问题；而策略不跟随速度命令、奖励曲线像随机游走，多半是\emph{MDP 层}的观测/奖励配置问题。第 2 步那套 MDP 逻辑的\emph{组织方式}，就是后文 \cref{sec:rlmc-eco-manager} 要讲的 Manager-Based 架构——它把这七步从一个巨型 \verb|step| 函数里拆成可独立替换的模块。
\end{insight}

\paragraph{角色三·调试仪器（ground-truth 读取）。}仿真相比真机有一个常被忽视的优势：可以读取\emph{完整}的物理状态。真机上你只能靠有噪声的部分传感器；仿真里你能直接读每个关节的精确 \verb|qpos|/\allowbreak\verb|qvel|、每个接触点的力、每个物体的质心位姿。这对调试奖励函数、诊断策略行为至关重要——\cref{ins:rlmc-eco-zerorandom} 的 zero/random 体检就是这一角色的最小用法。

\begin{center}
\small
\begin{tabular}{@{}lll@{}}
\toprule
调试能力 & 真机 & 仿真 \\
\midrule
关节角/速度 & 编码器（有噪声） & 精确 \verb|qpos|/\allowbreak\verb|qvel| \\
接触力 & 力传感器（若有） & 精确接触法向力/摩擦力 \\
质心位置 & 需状态估计器 & 直接读取 \\
地形信息 & 难以获取 & heightfield 精确数据 \\
状态保存/恢复 & 几乎不可能 & 可存完整状态并从任意点恢复 \\
\bottomrule
\end{tabular}
\end{center}

\subsection{三框架在「三角色」上的分工对比}

把三个角色当成尺子，三框架的定位差异立刻清晰（数值/版本以 2026 年中为快照，详见各小节核对）：

\begin{center}
\small
\begin{tabularx}{\linewidth}{@{}lLLL@{}}
\toprule
角色 & mjlab & Isaac Lab & UniLab \\
\midrule
物理世界 & MuJoCo Warp（凸优化接触） & PhysX / Newton(Beta) & MuJoCoUni / MotrixSim（\textbf{CPU} 多线程） \\
MDP 机器 & Manager-Based（借鉴 Isaac Lab） & Manager-Based（首创） & \textbf{非} Manager-Based（\verb|NpEnv| 子类） \\
调试仪器 & Viser（web 3D）+ 直接读 \verb|mjData| & Isaac Sim Viewer + RTX & Viser（web 3D）+ 直接读 \verb|NpEnvState| + NaN 守卫 \\
\bottomrule
\end{tabularx}
\end{center}

\begin{note}[关于第三个框架 UniLab：它\emph{不是}第三个 Manager-Based 框架]
本部按「三大框架对比」组织实战。前两家（mjlab / Isaac Lab）是「GPU 仿真主导」范式的两个 Manager-Based 框架；第三家 \textbf{UniLab}（论文 arXiv:2605.30313，副标题 \emph{A Heterogeneous Architecture for Robot RL Beyond GPU-Dominant Paradigms}）则\textbf{刻意正交}：它把物理仿真留在 \textbf{CPU}（MuJoCoUni / MotrixSim 多线程批量步进），把策略学习放在 \textbf{GPU}（PPO/APPO/SAC/TD3），二者经\emph{共享内存 IPC} 解耦（见 \cref{sec:rlmc-eco-manager}）。这意味着 UniLab \emph{不采用} Manager 体系——它没有 \Verb[breaklines]|ObservationManager|/\allowbreak\verb|RewardManager|/\allowbreak\verb|EventManager| 这些类，MDP 逻辑集中在 \verb|NpEnv| 子类的 \Verb[breaklines]|update_state()| 里。本书把 UniLab 当作「第三条工程路线」并入对比：凡 UniLab 与前两家\emph{设计同源}处直接并表，凡 UniLab \emph{独有}（异构 IPC、CPU 物理）或\emph{真实欠缺}（无 RTX 视觉、无 Manager 类）处都如实标注，不强行套用 Manager 词汇。本章正文凡 UniLab 列均为对源码与实跑核对后的真实信息（事实来源见 \cref{sec:rlmc-eco-tools} 延伸阅读）。
\end{note}

\begin{pitfall}[把「渲染漂亮」当成「适合训练」]
Isaac Sim 的 RTX 实时光追画面非常逼真，容易让人误以为它在所有任务上都更优。但对\emph{无视觉}的 locomotion 训练，渲染质量与训练吞吐、显存占用无关——你需要的是几千个并行环境快速 step，而非一帧逼真画面。逼真渲染只在\emph{视觉策略}（RGB/Depth 输入）任务里才是刚需。把这两件事混在一起，会让你为纯 locomotion 任务背上不必要的安装复杂度与显存开销（见 \cref{sec:rlmc-eco-select} 决策树）。
\end{pitfall}

\begin{practice}[练习·三角色]
\begin{enumerate}
  \item \textbf{[分析]} 你的机器人足端持续穿入地面、整机缓慢下沉。这个 bug 更可能出在「物理世界」还是「MDP 机器」角色？你会先检查哪一类参数？
  \item \textbf{[实操]} 在 mjlab 上对一个 velocity 任务跑 \Verb[breaklines]|--agent zero|，用调试仪器角色读出机器人此刻的若干关节角，判断默认姿态是否合理（提示：与官方 MJCF 的 keyframe 对照）。
  \item \textbf{[设计]} 列出三个仿真\emph{无法}完美模拟的物理现象（如柔性绳索、液体、软体接触），并思考这些任务是否仍可用 sim-to-real RL 范式解决。
\end{enumerate}
\end{practice}

% ════════════════════════════════════════════════════════════════
\section{为什么必须是 GPU 并行仿真}
\label{sec:rlmc-eco-gpu}

\subsection{动机：数据需求与真实世界成本的数量级鸿沟}

机器人 RL 的核心矛盾，可一句话概括：\textbf{策略学习需要海量试错数据，而真实世界的每一次试错都有物理代价}。这不是程度差异，而是数量级鸿沟。以四足 locomotion 为例，一个 PPO 策略从随机初始化到学会稳定行走，典型需要 $10^8$ 到 $10^9$ 个环境 step。下表把这个量级翻译成时间：

\begin{center}
\small
\begin{tabular}{@{}lll@{}}
\toprule
参数 & 值 & 说明 \\
\midrule
控制频率 & $50\,$Hz & 每秒 50 个 policy step \\
训练所需 step & $\sim 10^8$ & 典型 PPO 四足任务（单组超参） \\
单环境实时运行 & $10^8/50=2\times10^6\,$s $\approx 23\,$天 & 不停机、不维修 \\
4096 并行环境 & $2\times10^6/4096\approx 500\,$s $\approx 8\,$分钟 & GPU 并行的数量级优势 \\
\bottomrule
\end{tabular}
\end{center}

真实吞吐不会线性扩展（kernel 调度、reset、传感器模拟都有开销），但数量级差异已足够说明：\emph{GPU 并行仿真对机器人 RL 不是锦上添花，而是刚需}。这也呼应 \cref{ins:rlmc-eco-analogy} 的成本论——把 23 天压到 8 分钟，研究迭代才从「一天一组实验」变成「一小时十组」。

\subsection{CPU 多进程 vs GPU 并行：本质不是「核心多」}

GPU 并行不是凭空出现的。它对应一次计算范式的切换：

\begin{center}
\small
\begin{tabular}{@{}lll@{}}
\toprule
维度 & CPU 多进程（2021 前主流） & GPU 并行（2021 后） \\
\midrule
并行度 & 16--64（受核心数限制） & 4096--65536（受显存限制） \\
数据位置 & 各进程独立内存 & 连续 GPU 显存 \\
通信开销 & 进程间通信（微秒级） & 无通信（同一显存空间） \\
与 RL 算法交互 & 需 CPU$\to$GPU 数据搬运 & 数据已在 GPU 上，零搬运 \\
编程模型 & 标准 Python 多进程 & CUDA kernel / Warp \\
\bottomrule
\end{tabular}
\end{center}

\begin{insight}[GPU 并行的核心优势是「数据零搬运」，不只是「核心多」]\label{ins:rlmc-eco-zerocopy}
端到端看一次训练迭代：\verb|env.step()| 产出 obs/reward $\to$ 策略网络前向算 action $\to$ rollout buffer 存储 $\to$ PPO 反向更新。在 GPU 并行范式下，这一整条流水线\emph{从头到尾都在 GPU 显存上}，没有任何 CPU--GPU 搬运。而在 CPU 多进程范式下，即使你有\emph{无限多}的 CPU 核心，每次 rollout 后仍要把数据从 CPU 搬到 GPU 做策略前向、再搬回 CPU 做下一次 \verb|env.step()|——这个搬运延迟（微秒级）在 4096 环境规模下累积成无法消除的瓶颈。所以「为什么 CPU 加核也追不上 GPU」的答案，是搬运，不是算力。
\end{insight}

\begin{insight}[多视角·有人偏要把物理放回 CPU：UniLab 的异构反范式]\label{ins:rlmc-eco-heterogeneous}
上一条把 CPU 多进程说成「追不上 GPU」，但这只是\emph{硬把整条管线钉在一种设备上}时的结论。第三个框架 \textbf{UniLab} 给出了反例：它\emph{故意}把物理仿真留在 CPU（MuJoCoUni / MotrixSim 多线程批量步进），只把策略学习放在 GPU，两端用\textbf{共享内存 IPC} 解耦——transition 单向从 collector（CPU 子进程）流向 learner（GPU 父进程），actor 权重经 \Verb[breaklines]|SharedWeightSync|（版本化扁平 buffer + \verb|int64| 版本号、单写多读、零序列化）单向流回。它\emph{没有}消除 CPU--GPU 边界，而是把这条边界做成\textbf{显式、可计量、可预取}的 H2D 管线（pinned-host + 双缓冲 + 可选原生扩展），让「CPU 采样打包」与「GPU 训练」流水线重叠。这样换来三件 GPU-sim 给不了的东西：(1) 不依赖任何 GPU 物理后端——任何能跑 MuJoCo 的机器都行；(2) collector 不必等 learner，异步 off-policy（APPO/SAC）成了一等公民；(3) 接触求解在 CPU 全精度。代价也实在：放弃了 CUDA-graph 那种把整条物理步固化的极致单卡并行，并多出一个新失败面——\verb|/dev/shm| 共享内存预算（UniLab 专门有 \Verb[breaklines]|memory_budget.py| 在分配前估算、超 80\% 直接 \verb|MemoryError|）。所以「GPU 并行是不是唯一答案」的诚实回答是：\emph{当 GPU-sim 后端可用且追求单卡极致吞吐时是；当 CPU 核多、无强 GPU-sim 后端、或要异步解耦时，把物理留在 CPU 反而是更优的工程选择}。记住这条多视角，你才不会把「GPU 并行」误当成机器人 RL 的唯一信仰。
\end{insight}

GPU 的执行模型是 SIMT（Single Instruction, Multiple Threads）：同一条指令同时在数千线程上执行，每线程操作不同数据。这天然契合环境并行——每个环境的 \verb|step()| 跑\emph{完全相同}的计算流程（同一指令），但各自操作\emph{独立}的状态数据（不同数据）。这种「相同计算、不同数据」的模式在 GPU 上极高效：指令缓存命中率高、状态数据连续排布可合并访存（memory coalescing）、环境间无通信需求。

\begin{insight}[SIMT 直接约束你怎么写环境]\label{ins:rlmc-eco-simt}
理解 SIMT 不是为了写 CUDA，而是因为它\emph{约束}了环境设计。两类做法在 GPU 并行里低效或不允许：(1) \Verb[breaklines]|if env_id == 0: do_special()| 这类按环境分支会导致 warp 发散（warp divergence），GPU 效率骤降——正确替代是用 mask 张量：\Verb[breaklines]|result = mask*special + (1-mask)*normal|；(2) 不同环境用不同 reward 函数违反「同一指令」原则——正确做法是所有环境算相同 reward，用权重 mask 控制哪些项生效。好在 Manager-Based 架构已替你处理了大部分：你在 config 里写的 \verb|ObsTerm|/\allowbreak\verb|RewTerm|/\allowbreak\verb|EventTerm| 都被框架自动转成符合 SIMT 的张量操作，你无需手写 kernel。这也解释了\emph{为什么}框架要做成 Manager-Based 的——它是 SIMT 约束下「让研究者像搭积木」的工程答案。
\end{insight}

\subsection{数据布局与 action repeat：两个易被忽视的工程细节}

\paragraph{SoA 而非 AoS。}对 $N$ 个环境、每个 $D$ 维状态，GPU 偏好 SoA（Structure of Arrays）布局：同一维度的数据在显存里连续，当所有线程同时读同一维度时可合并为一次访存。这就是为什么框架里的张量形状几乎总是 \Verb[breaklines]|(num_envs, dim)| 而非 \Verb[breaklines]|(dim, num_envs)|——这不是随意约定，而是访存模式的直接要求。mjlab 与 Isaac Lab 都用 SoA。

\paragraph{action repeat（decimation）。}策略决策频率（如 $50\,$Hz）通常低于物理步频（如 $200\,$Hz），于是每次 \verb|env.step()| 内部调用多次 \verb|sim.step()|：

\begin{center}
\small
\begin{tabular}{@{}lll@{}}
\toprule
参数 & 含义 & 典型值 \\
\midrule
\verb|dt| & 物理时间步长 & $0.005\,$s（$200\,$Hz） \\
\verb|decimation| & 每个 policy step 的物理子步数 & 4 \\
\verb|policy_dt| & 决策周期（\verb|dt| 乘 \verb|decimation|） & $0.02\,$s（$50\,$Hz） \\
\bottomrule
\end{tabular}
\end{center}

\begin{insight}[action repeat 是「决策频率」与「物理精度」之间的桥]\label{ins:rlmc-eco-decimation}
策略不必每个物理步都决策——就像人走路时大脑不以 $200\,$Hz 控制每块肌肉，而是低频发高层指令，底层反射回路做高频调节。decimation 太小（如 1）使策略输出变化过于细碎、训练负担重；太大（如 20）使策略反应迟钝、来不及修正误差；locomotion 常用 4。这个参数会贯穿后文的延迟建模——观测延迟的 lag 单位正是 env step 而非 physics step（详见 \cref{ch:rlmc-obsact} 的 delay 阶段）。
\end{insight}

\subsection{吞吐量、CUDA Graph 与显存上限}

不同框架在相同硬件上的吞吐差异，比想象中小。下表是\emph{量级参考}（\pz{待核实：具体 steps/s 强依赖任务与硬件，源稿给出的 A100 上 mjlab/Isaac Lab 约 $4$--$5\times10^5$ steps/s 为社区量级估计，未逐一复现}）：

\begin{center}
\small
\begin{tabular}{@{}lllll@{}}
\toprule
框架 & 物理引擎 & 并行度 & 量级（steps/s） & 备注 \\
\midrule
CPU MuJoCo & MuJoCo & 1 & $\sim 10^3$ & 单核 baseline \\
mjlab & MuJoCo Warp & 4096 & $\sim 10^5$ & \rebuilt{量级参考} \\
Isaac Lab & PhysX & 4096 & $\sim 10^5$ & \rebuilt{量级参考} \\
MuJoCo Playground & MJX(JAX) & 4096 & $\sim 10^5$ & \rebuilt{量级参考} \\
\bottomrule
\end{tabular}
\end{center}

要点：GPU 并行相比 CPU 单核有\emph{两到三个数量级}的提升（这与 Isaac Gym 原论文「2--3 orders of magnitude」一致\,\cite{makoviychuk2021isaacgym}），但 mjlab 与 Isaac Lab 之间在同一数量级、差异不大——\textbf{吞吐量不是选框架的主要依据}，物理精度、API 设计、生态支持更重要。

\begin{insight}[别被「首跑小 env 数的低吞吐」吓到，且学会自己测]\label{ins:rlmc-eco-throughput}
上表的 $\sim 10^5$ steps/s 是 \textbf{4096 并行环境} 量级的稳态峰值，\emph{不是}你第一次小规模冒烟时会看到的数字。在 RTX 4090 上实测 mjlab \Verb[breaklines]|Velocity-Flat-Unitree-G1| 仅 \Verb[breaklines]|--env.scene.num-envs 64| 时，稳态约 \textbf{$1.6$--$2.5\times10^3$ steps/s}——比峰值低了近两个数量级。原因有二：(1) 64 个环境远填不满 GPU 的数千个 SIMT 通道（kernel 未饱和），吞吐受启动/调度而非算力主导；(2) \textbf{首个}迭代还要付一次性的 JIT 编译与 CUDA-graph 录制开销，所以\emph{第一行 steps/s 总是最低}，跑几十迭代后才升到稳态。结论：冒烟用小 env 数看到几千 steps/s 是\emph{正常}的，别误判成「装坏了」；要看真实吞吐，把 \Verb[breaklines]|--env.scene.num-envs| 加到 $2048$/$4096$ 再读稳态行。\textbf{怎么自己测}——不必写计时代码，三框架的训练循环都\emph{自带}每迭代打印 \verb|steps/s|（mjlab/Isaac Lab 经 RSL-RL 的 \Verb[breaklines]|OnPolicyRunner|，UniLab 同理）：跑十几迭代、读\emph{稳定后}那几行的数字，就是你这台机器在该 env 数下的实测吞吐。把这个数记下来当 baseline，日后某次改动让它掉一半，往往就是无意中触发了 CUDA-graph 重建或把某段计算挪回了 CPU（呼应 \cref{ins:rlmc-eco-zerocopy}）。
\end{insight}

即便数据全在 GPU，每次 \verb|env.step()| 仍要启动多个 CUDA kernel（接触检测、力计算、积分、obs 拼接、reward……），每个 kernel 启动都有微秒级 CPU 端开销。\textbf{CUDA Graph} 把一串 kernel 启动序列「录制」下来，之后一次 CPU 调用即可重放整段——mjlab 与 Isaac Lab 都支持。但它有硬限制：\emph{计算图必须静态}，不允许动态分支。框架常用「延迟 reset」（记录待 reset 的环境 ID，待所有环境 step 完成后统一 reset）来规避「不同环境在不同时刻 reset」造成的分支。

\begin{insight}[GPU 并行的上限是显存，不是核心数]\label{ins:rlmc-eco-vram}
每个并行环境要存机器人状态（\verb|qpos|/\allowbreak\verb|qvel|）、接触缓冲、obs/action/reward 张量。一个粗略但够用的估算公式是
\begin{equation}
\texttt{max\_envs}\ \approx\ \frac{(\texttt{VRAM\_GB}-2)\times 1000}{\texttt{per\_env\_MB}},
\label{eq:rlmc-eco-maxenvs}
\end{equation}
其中预留 $2\,$GB 给 PyTorch 框架与策略网络。它不精确，却足以在训练\emph{前}判断「4096 环境在我的卡上跑不跑得动」。一个关键提醒：当任务含视觉传感器（RGB/Depth）时，每环境显存会暴涨 $10$--$50$ 倍——这正是视觉策略任务往往需要 $80\,$GB 卡的原因，也是 \cref{sec:rlmc-eco-select} 决策树把「是否需要视觉」放在第一问的工程根据。
\end{insight}

\begin{pitfall}[三个关于 GPU 并行的常见误解]
\textbf{一·以为 GPU 并行\emph{总}更快}：对极简任务（如 CartPole），kernel 启动开销可能超过计算本身，此时 CPU 多进程反而更快。GPU 的优势在高自由度机器人（四足/人形）上才真正体现。\textbf{二·忽视 CUDA Graph 的静态图限制}：若环境有复杂动态行为（地形切换、curriculum 升级），CUDA Graph 可能失效，表现为「某些 step 突然变慢」。\textbf{三·混淆「环境数」与「有效样本量」}：4096 个\emph{初始条件相同、无域随机化}的环境产生高度相关的数据，对策略更新的贡献远小于 4096 个独立样本——数量不等于多样性。
\end{pitfall}

\begin{practice}[练习·GPU 并行]
\begin{enumerate}
  \item \textbf{[推理]} 解释为什么 CPU 多进程范式即便有 4096 个核心，训练速度仍追不上 GPU 并行（提示：\cref{ins:rlmc-eco-zerocopy} 的搬运延迟）。
  \item \textbf{[计算]} 查出你常用 GPU 的显存，用 \cref{eq:rlmc-eco-maxenvs} 估算每环境约 $2\,$MB 状态时最多能并行多少个人形环境；再估算每环境含相机约 $50\,$MB 时的数字。
  \item \textbf{[计算]} 任务用 \verb|dt=0.005|（$200\,$Hz）、\verb|decimation=4|。决策频率是多少 Hz？若改 \verb|decimation=8|，频率变多少？这对训练效果可能有何影响？
\end{enumerate}
\end{practice}

% ════════════════════════════════════════════════════════════════
\section{生态演进：每一代解决了什么痛点}
\label{sec:rlmc-eco-evolution}

\subsection{为什么要懂演进脉络}

2026 年的 GitHub 上至少有六个活跃的机器人 RL 仿真框架。一个博士新生面对它们会困惑：「它们什么关系？我该用哪个？」要回答，必须理解演进脉络——\emph{每一代框架都是为解决上一代的痛点而生}。不懂这条线，你会看不懂为什么 2023 年的论文用 Isaac Gym 而 2026 年的用 mjlab，可能花几天配好一个已停止维护的框架，甚至在不适合你任务的框架上耗费数月。

\subsection{六个阶段：核心创新与遗留痛点}

\paragraph{阶段一·CPU MuJoCo（2012--2021）。}MuJoCo（Multi-Joint dynamics with Contact）由 Todorov、Erez、Tassa 在 IROS 2012 发表\,\cite{todorov2012mujoco}。经核对，它于 2021 年 10 月被 Google DeepMind 收购，2022 年 5 月以 Apache 2.0 完整开源\,\cite{todorov2012mujoco}。其核心创新是\textbf{广义坐标 + 凸优化接触}：接触约束不像传统 LCP 那样求解互补性问题（病态时可能发散），而是松弛为凸优化——\emph{求解总能成功}。这使它在足式机器人（接触频繁建立/断开）与灵巧手（手指与物体复杂接触）中表现出色。遗留痛点：\textbf{单线程 CPU，吞吐低}——训练四足要数小时到数天。

\begin{paper}[MuJoCo：广义坐标 + 凸优化接触的物理引擎]\label{paper:rlmc-eco-mujoco}
\textbf{问题}：如何在接触密集的多关节系统中兼顾仿真精度与速度，且求解稳定。\textbf{核心思想}：用广义（关节）坐标表达动力学，把硬接触的互补性约束松弛为\emph{凸优化}问题求解。\textbf{关键结果}：凸优化保证全局最优解、求解总能收敛，代价是允许极微小穿透与微滑（slip）；在足式与灵巧操作上取得当时最佳的精度-速度平衡\,\cite{todorov2012mujoco}。\textbf{与本书关系}：mjlab 的 MuJoCo Warp 本质是「保留这套接触精度，但搬到 GPU 上跑」——理解 \cref{paper:rlmc-eco-mujoco} 才懂 \cref{sec:rlmc-eco-manager} 里 mjlab「回到 MuJoCo 物理」的动机。\textbf{局限}：原版单线程 CPU，吞吐量不足以支撑大规模并行 RL。
\end{paper}

接触模型的差异值得先建立直觉，因为它是后续物理引擎选型的核心。传统游戏引擎（PhysX/Bullet/ODE）用\textbf{硬接触}：求解一组力使穿透量为零，这是 LCP（线性互补问题），数学上 NP-hard，求解器只能近似（如 PhysX 的 TGS 迭代），接触点多或几何复杂时可能不收敛。MuJoCo 用\textbf{软接触}：允许微小穿透，用凸优化算使穿透最小的接触力，求解\emph{总能成功}，代价是接触面有极微小穿透与微滑。

\begin{pitfall}[接触参数选错的三种典型症状]
MuJoCo 用 \verb|solref|（接触刚度/阻尼）与 \verb|solimp|（约束阻抗）控制接触行为。选错的三种症状：\textbf{穿透}（\verb|solref| 太软）——足端穿过地面，机器人「沉」入地下；\textbf{弹跳}（\verb|solref| 太硬 + 阻尼不足）——物体在接触面反复振荡；\textbf{微滑}（软接触固有特性）——静止物体在斜面缓慢下滑，可调 \verb|impratio| 缓解。对多数 locomotion 任务，MuJoCo 默认值已够好，只在接触行为明显异常时才手动调。\pz{源稿给出各材质 solref 参考值（如金属 [0.005,1.0]）引自 ROBOLAWEB cheat sheet，本书未独立复现，使用前请对照官方文档。}
\end{pitfall}

\paragraph{阶段二·Isaac Gym（2021--2024）。}NVIDIA 在 NeurIPS 2021 发布 Isaac Gym\,\cite{makoviychuk2021isaacgym}，\textbf{首次实现 GPU 并行仿真 + 端到端 GPU 训练}，是机器人 RL 工程化的分水岭。后端为 PhysX。它对社区影响巨大：ETH RSL 的 \verb|legged_gym|（Rudin 等，CoRL 2021\,\cite{rudin2021minutes}）在其上确立了四足 locomotion 的标准训练范式，后续几乎所有足式 RL 论文都以它为基。但 Isaac Gym 的\textbf{单体架构}痛点很快暴露：一个任务文件常超 2000 行，obs/reward/reset/域随机化全耦合在一个巨型类里，改任一组件都可能引入隐蔽 bug。

\begin{paper}[Isaac Gym / legged\_gym：GPU 并行 locomotion 的奠基]\label{paper:rlmc-eco-isaacgym}
\textbf{问题}：CPU 仿真吞吐不足，四足 locomotion 训练动辄数小时到数天。\textbf{核心思想}：把物理仿真与神经网络训练\emph{都}放在 GPU 上，数据从物理 buffer 直接传给 PyTorch 张量，不经 CPU（\cref{ins:rlmc-eco-zerocopy} 的「零搬运」即源于此）\,\cite{makoviychuk2021isaacgym}；\verb|legged_gym| 在此之上给出大规模并行 PPO 的 locomotion 范式\,\cite{rudin2021minutes}。\textbf{关键结果}：相较 CPU 仿真 + GPU 训练有 2--3 个数量级加速\,\cite{makoviychuk2021isaacgym}；平地四足策略可在数分钟内训成，崎岖地形约二十分钟\,\cite{rudin2021minutes}。\textbf{与本书关系}：理解它的\emph{单体架构痛点}，才理解 Manager-Based 架构（\cref{sec:rlmc-eco-manager}）为何不可避免。\textbf{局限}：单体类不可组合、接触精度不如 MuJoCo、仅支持 PhysX；Isaac Gym 本身已停止维护。
\end{paper}

\begin{codebox}[Python]{单体架构（legged\_gym 风格，简化示意）:一个类装下所有逻辑}
class LeggedRobot(BaseTask):
    def step(self, actions):
        self.pre_physics_step(actions)   # action 处理
        self.gym.simulate(self.sim)      # 物理 step
        self.post_physics_step()         # obs/reward/reset 全混在这里

    def compute_observations(self):      # 200+ 行 obs 计算
        self.obs_buf = torch.cat([...], dim=-1)

    def compute_reward(self):            # 300+ 行 reward 计算
        self.rew_buf = ...

    def reset_idx(self, env_ids):        # 150+ 行 reset 逻辑
        ...
    def _push_robots(self):              # 域随机化也在同一个类里
        ...
\end{codebox}

\paragraph{阶段三·Isaac Lab（2024--）。}NVIDIA 2024 年发布 Isaac Lab\,\cite{mittal2025isaaclab} 作为 Isaac Gym 继任者，核心创新是 \textbf{Manager-Based 架构}——把 MDP 的每个组件（obs/action/reward/termination/event/command/curriculum）拆到独立 Manager，实现关注点分离。后端默认 PhysX，3.0 Beta 起可经 Newton 用 MuJoCo Warp。优势：架构清晰、多物理后端、RTX 渲染（视觉任务）、社区最大。痛点：\textbf{安装复杂}（需 Isaac Sim/Omniverse 或 kit-less 配置）、启动延迟较大、对 MuJoCo 物理的支持仍在 Beta。

\paragraph{阶段四·MuJoCo Playground / MJX（2024--2025）。}Google DeepMind 推出 MuJoCo 的 JAX 加速版 MJX 及其上的 MuJoCo Playground\,\cite{zakka2025playground}，用 JAX \verb|vmap| 实现批量仿真，物理精度与 CPU MuJoCo 一致，且\emph{可微分}。其独特能力是基于 Madrona 的\emph{批量渲染}——可在 GPU 上同时为数千环境生成 RGB/Depth\,\cite{zakka2025playground}。痛点：JAX 生态与主流 PyTorch RL 算法库（RSL-RL/RL Games）不直接兼容。

\paragraph{阶段五·Newton（2025--）。}Newton 由 NVIDIA、Google DeepMind、Disney Research 联合发起，2025 年起由 Linux Foundation 治理为开源项目（Apache 2.0）\,\cite{newton2025}。目标是\textbf{统一多个物理求解器}在一个 API 下，让用户在同一任务中切换后端而不改 MDP 代码。其中 Disney 的 \textbf{Kamino} 求解器专攻\emph{闭环机构}（平行连杆腿、灵巧手）——多数引擎难以稳定仿真的拓扑\,\cite{newton2025}。

\begin{paper}[Newton：多求解器统一的开源物理引擎]\label{paper:rlmc-eco-newton}
\textbf{问题}：不同任务（locomotion / 闭环机构 / 软体 / 颗粒材料）需要不同物理求解器，过去要维护多套环境代码才能跨引擎对比。\textbf{核心思想}：在 NVIDIA Warp + OpenUSD 之上，用统一 API 容纳多个求解器（含 MuJoCo Warp 主求解器与 Disney 的 Kamino 闭环求解器），切换后端只改配置\,\cite{newton2025}。\textbf{关键结果}：据 NVIDIA 技术博客，MuJoCo Warp 相对 MJX 在 RTX PRO 6000 Blackwell 上 locomotion 约 $252\times$、manipulation 约 $475\times$\,\cite{newton2025blog}。\textbf{与本书关系}：Newton 正让 Isaac Lab 与 mjlab 在\emph{物理后端}层面趋同（两者都可用 MuJoCo Warp），未来选框架更多取决于 API/生态而非引擎——这是本书选「双框架对比」教学的核心理由。\textbf{局限}：Isaac Lab 对 Newton 的集成仍 Beta。
\end{paper}

\begin{pitfall}[别盲信「up to」峰值 benchmark]
\cref{paper:rlmc-eco-newton} 引的 $252\times$/$475\times$ 是 NVIDIA 官方在特定硬件（RTX PRO 6000 Blackwell）、特定任务（locomotion / manipulation）上对 MJX 的\emph{峰值}（up to），非平均值；换 RTX 4090 数字降到约 $152\times$/$313\times$\,\cite{newton2025blog}。教学选型里，性能差异\emph{不}是主要依据。经核对 NVIDIA 技术博客，「Newton Beta 作 Isaac Lab 后端时 in-hand 灵巧操作比 PhysX 快至多 $65\%$」一句\emph{属实}\,\cite{newton2025blog}。\pz{源稿曾写「Newton 含 7 个求解器」，经核对官方 \texttt{newton.\allowbreak solvers} API 实为 8 个具体求解器类——\texttt{SolverMuJoCo}、\texttt{SolverXPBD}、\texttt{SolverVBD}、\texttt{SolverImplicitMPM}、\texttt{SolverFeatherstone}、\texttt{SolverSemiImplicit}、\texttt{SolverStyle3D}、\texttt{SolverKamino}（另有抽象基类 \texttt{SolverBase}）；具体个数随版本增删，以你所装 Newton 的 \texttt{newton.\allowbreak solvers} 为准。}
\end{pitfall}

\paragraph{阶段六·mjlab（2026--）。}UC Berkeley 的 Zakka、Liao、Yi、Le Lay、Sreenath、Abbeel 团队 2026 年发布 mjlab\,\cite{mjlab2026}。设计哲学一句话：\textbf{Isaac Lab 的 Manager-Based API + MuJoCo Warp 的 GPU 物理 + 一条命令安装}。后端为 MuJoCo Warp（Google DeepMind 与 NVIDIA 共同维护的 MuJoCo GPU 加速版）。优势：轻量、MuJoCo 原生精度、Isaac Lab 风格 API、可直接访问 \verb|MjModel|/\allowbreak\verb|MjData|。痛点：社区规模与内置任务数不如 Isaac Lab、不支持 RTX 渲染、仅支持 RSL-RL（PPO）。

\begin{paper}[mjlab：轻量级 GPU 机器人学习框架]\label{paper:rlmc-eco-mjlab}
\textbf{问题}：Isaac Lab 的 Manager-Based API 很好，但 Isaac Sim 依赖让安装与启动笨重；MuJoCo 物理精度高，却缺一个 PyTorch 生态下的轻量框架封装。\textbf{核心思想}：把 Isaac Lab 的 Manager-Based API 直接搭在 MuJoCo Warp 上，单条 \Verb[breaklines]|pip install| 即可，并\emph{直接暴露} MuJoCo 原生数据结构\,\cite{mjlab2026}。\textbf{关键结果}：随附 velocity tracking / motion imitation / manipulation 等参考任务；学会 Isaac Lab 后迁移到 mjlab 认知成本低（核心 Manager 概念一一对应）\,\cite{mjlab2026}。\textbf{与本书关系}：mjlab 是本部「轻量/接触密集」实战的主线框架之一。\textbf{局限}：仅 RSL-RL 后端、无 RTX 渲染。\pz{源稿引用的 mjlab arXiv 号 2601.22074 与版本/Stars 等细节中，arXiv 号与作者经核对属实；Stars 等时效数据本书不固定。}
\end{paper}

\begin{insight}[一条贯穿六阶段的模式：保留核心优势、消除核心痛点]\label{ins:rlmc-eco-pattern}
六个阶段不是孤立的，而是一条「继承-改良」链：Isaac Gym 保留了 GPU 并行、消除了 CPU 吞吐痛点；Isaac Lab 保留 GPU 并行、用 Manager 架构消除单体不可组合痛点；mjlab 保留 Manager API、用 MuJoCo Warp + pip 消除「物理精度不足 + 安装笨重」痛点；Newton 则保留各家引擎优势、用统一 API 消除「跨引擎要维护多套代码」痛点。看懂这条模式，你读任何\emph{新}框架时都能一句话定位它：「它保留了谁的什么，又想消除谁的什么痛点。」
\end{insight}

\subsection{演进时间线（速查）}

\begin{center}
\small
\begin{tabular}{@{}llll@{}}
\toprule
年份 & 框架/项目 & 后端 & 解决的痛点 / 留下的痛点 \\
\midrule
2012 & MuJoCo（Todorov） & CPU 广义坐标 & 接触精度高 / 单线程吞吐低 \\
2021 & Isaac Gym（NVIDIA） & PhysX(GPU) & GPU 并行 / 单体架构 \\
2021 & legged\_gym（ETH RSL） & PhysX(GPU) & locomotion 标准范式 / 不可组合 \\
2024 & Isaac Lab（NVIDIA） & PhysX/Newton & Manager 架构 / 安装复杂 \\
2024 & MuJoCo Playground/MJX & MJX(JAX) & 可微分 + 批量渲染 / JAX 生态 \\
2025 & Newton（LF 治理） & 多求解器 & 统一引擎 / Isaac Lab 集成 Beta \\
2026 & mjlab（UC Berkeley） & MuJoCo Warp & MuJoCo 精度 + pip 安装 / 仅 PPO \\
\bottomrule
\end{tabular}
\end{center}

\begin{pitfall}[三个关于生态的常见误解]
\textbf{一·以为 Isaac Gym 还是主流}：截至 2026 年它已被 Isaac Lab 与 mjlab 取代，新项目不应再基于它；但很多高星老项目仍用它，读其代码要理解架构限制。\textbf{二·把 MJX 与 MuJoCo Warp 当同一物}：MJX 基于 JAX（用于 MuJoCo Playground），MuJoCo Warp 基于 NVIDIA Warp（用于 mjlab，由 DeepMind + NVIDIA 共同维护）；物理精度相近，编程模型完全不同。\textbf{三·混淆「框架」与「物理引擎」}：Isaac Lab 是\emph{框架}、PhysX/MuJoCo Warp 是\emph{物理引擎}；同一引擎可被不同框架使用（MuJoCo Warp 被 mjlab 直接用、被 Isaac Lab 经 Newton 间接用）。
\end{pitfall}

\begin{practice}[练习·生态演进]
\begin{enumerate}
  \item \textbf{[梳理]} 画一张「框架-引擎」对应图，标明 mjlab、Isaac Lab、legged\_gym、MuJoCo Playground 各用什么物理引擎。
  \item \textbf{[对比]} 找一篇 2024 年与一篇 2026 年的足式机器人论文，对比其仿真框架，分析迁移的技术原因。
  \item \textbf{[应用]} 用 \cref{ins:rlmc-eco-pattern} 的「保留-消除」句式，给一个你听说过但本章没讲的框架（如 Genesis）做一句话定位。
\end{enumerate}
\end{practice}

% ════════════════════════════════════════════════════════════════
\section{Manager-Based 架构：三框架对比实战}
\label{sec:rlmc-eco-manager}

\subsection{模块功能与在管线中的位置}

\cref{ins:rlmc-eco-simenv} 指出 \verb|env.step()| 的七步 MDP 逻辑需要某种\emph{组织方式}。Manager-Based 架构就是这个组织方式：把 MDP 的每个维度交给一个独立 Manager 管理，每个 Manager 由若干\emph{项}（term）组成，项之间通过 config 组合。这把复杂任务从一个巨型 \verb|step| 函数里解耦出来——\emph{改一个奖励不再牵动观测}。这套架构 Isaac Lab 首创、mjlab 借鉴；UniLab 则\emph{不走} Manager 路线，但要处理的那些 MDP 维度一个不少，只是组织在别处（下一小节逐项对齐）。

经核对 Isaac Lab 官方文档\,\cite{mittal2025isaaclab}，其公开 Manager 覆盖 observation / action / reward / termination / command / curriculum / event / recorder 等维度。下表给出职责与「你何时会改它」，并标注三框架支持情况：

\begin{center}
\small
\begin{tabular}{@{}llll@{}}
\toprule
Manager & 职责 / 何时修改 & mjlab & Isaac Lab \\
\midrule
ObservationManager & 算 obs 并拼张量 / 改策略输入 & \checkmark & \checkmark \\
ActionManager & 处理 action 并下发执行器 / 改动作空间 & \checkmark & \checkmark \\
RewardManager & 算各 reward 项并加权 / 改奖励 & \checkmark & \checkmark \\
TerminationManager & 判 episode 结束 / 改终止条件 & \checkmark & \checkmark \\
EventManager & 域随机化与事件 / 改 DR 策略 & \checkmark & \checkmark \\
CommandManager & 生成指令并管生命周期 / 改指令范围 & \checkmark & \checkmark \\
CurriculumManager & 按进度调难度 / 改课程 & \checkmark & \checkmark \\
RecorderManager & 录视频/状态 / 加录制 & \checkmark & \checkmark \\
\bottomrule
\end{tabular}
\end{center}

\begin{note}[一个易传错的「Manager 数目」说法]
源稿称「Isaac Lab 八大 Manager、mjlab 第九个 MetricsManager 独有」。经在服务器实装的 mjlab 1.4.0 上核对源码：\Verb[breaklines]|mjlab.managers| 确实多出一个 \textbf{MetricsManager}（Isaac Lab 无对应——它无 weight、不 $\times$dt、不除 episode 长度，专报「真·逐步回合平均」），这一「mjlab 独有」说法\emph{成立}；mjlab 另有 Action/Command/Curriculum/Event/Observation/Recorder/Reward/Termination 共九个管理器。但「Isaac Lab 恰好八个公开 Manager」这类\emph{精确计数}随版本变化、不同文档口径不一。\pz{待核实：两框架公开 Manager 的确切个数随版本浮动，请以各自当时版本的 API 文档为准；本书只依赖「每个 Manager 对应 MDP 一个维度、改一个不影响其他」这一稳定事实。}
\end{note}

上表只列了两家，因为\textbf{UniLab 没有 Manager 类}——但 Manager 要管的那些 MDP 维度，UniLab 一样要处理，只是\emph{散布在别处}。把「每个 Manager 的职责落在 UniLab 的哪里」对齐一遍，恰好能反向印证「Manager 把什么聚合成了什么」：

\begin{center}
\small
\begin{tabular}{@{}p{3.6cm}p{10.0cm}@{}}
\toprule
Manager 职责（mjlab/Isaac Lab） & UniLab 把它放在哪（经源码核对） \\
\midrule
Observation / Reward / Termination & 合并进 \Verb[breaklines]|NpEnv.update_state()| \emph{一处}算：obs 在 \verb|_compute_obs| 拼，reward 经 \Verb[breaklines]|run_reward_dispatch| 派发，terminated 直接判 \\
Reward 的「项+权重」 & 不是 \verb|RewTerm| 类，而是 Hydra \verb|reward.scales| 标量字典 + env 内 \verb|_reward_fns| 派发表（\verb|RewardContext| 传上下文） \\
Event（域随机化） & 独立 \Verb[breaklines]|DomainRandomizationManager|（init / reset / interval 三模式，对应 \verb|EventMode|） \\
Curriculum & 独立 \Verb[breaklines]|PenaltyCurriculum|（按 episode 长度缩放负权重）+ 地形课程 \\
Command & \verb|Commands|（速度命令采样 + heading 反馈） \\
Action & 无 \verb|ActionCfg| 类：\verb|apply_action| 里 \Verb[breaklines]|ctrl = act*action_scale + default_angles|（位置式 PD 残差） \\
Recorder / Metrics & 无对应 Manager 类；诊断靠 \verb|NanGuard|（自动 dump + \Verb[breaklines]|unilab-viz-nan| 离线复现）+ chrome-trace profiling \\
\bottomrule
\end{tabular}
\end{center}

\begin{insight}[Manager-Based 不是「唯一正确」，是一种\emph{聚合粒度}的选择]\label{ins:rlmc-eco-manager-granularity}
对照上表能看出一个常被忽略的事实：Manager-Based 的本质是\textbf{把 MDP 沿「维度」切开、每维一个可独立替换的盒子}。UniLab 把 obs/reward/termination \emph{合在一处}（\verb|update_state|）、却把 DR 和课程\emph{拆成独立 manager}——它的切分线和 Isaac Lab 不同，但同样在「关注点分离」与「一处看全」之间做权衡。所以「要不要 Manager」不是对错题：Manager 把消融做廉价（改一行 config）、代价是逻辑被打散到多个文件、新手要先建立「哪个 term 在哪个 manager」的地图；\verb|update_state| 一处算齐则一眼看全数据流、代价是想做消融得改函数体。看清这一点，你读任何框架时都能问出更准的问题：「它把 MDP 切成了几块、沿什么线切、谁能独立换」——而不是只会问「它有没有 RewardManager」。
\end{insight}

\subsection{配置详解：Manager-Based 最小骨架（四维配置表）}

把抽象落到代码。Manager-Based 任务的核心是一个 \verb|@configclass| 配置树，下面是「观测 + 奖励 + 事件」三个 Manager 的最小骨架（Isaac Lab 风格；mjlab 结构几乎同形，差异见 \cref{sec:rlmc-eco-manager-compare}）。先看代码，再用四维表解读关键配置项。

\begin{codebox}[Python]{Manager-Based 最小骨架（Isaac Lab 风格，简化示意）}
from isaaclab.utils import configclass
from isaaclab.managers import ObservationGroupCfg as ObsGroup
from isaaclab.managers import ObservationTermCfg as ObsTerm
from isaaclab.managers import RewardTermCfg as RewTerm
from isaaclab.managers import EventTermCfg as EventTerm
import isaaclab.envs.mdp as mdp

@configclass
class ObservationsCfg:
    @configclass
    class PolicyCfg(ObsGroup):                       # 部署可得信号（actor 用）
        joint_pos = ObsTerm(func=mdp.joint_pos_rel)  # 相对默认姿态的关节角
        joint_vel = ObsTerm(func=mdp.joint_vel_rel)
        base_ang_vel = ObsTerm(func=mdp.base_ang_vel)
    policy: PolicyCfg = PolicyCfg()

@configclass
class RewardsCfg:
    track_lin_vel = RewTerm(func=mdp.track_lin_vel_xy_exp, weight=1.5)
    track_ang_vel = RewTerm(func=mdp.track_ang_vel_z_exp, weight=0.75)

@configclass
class EventCfg:
    # mode 决定何时触发:startup（仅一次）/ reset（每次重置）/ interval（周期）
    randomize_mass = EventTerm(func=mdp.randomize_rigid_body_mass, mode="startup")
    push_robot = EventTerm(func=mdp.push_by_setting_velocity, mode="interval")
\end{codebox}

\begin{center}
\small
\begin{tabular}{@{}lp{2.6cm}p{3.0cm}p{3.4cm}@{}}
\toprule
配置项 & 默认值来源 & 修改影响 & 合理范围 / 与其他配置耦合 \\
\midrule
\Verb[breaklines]|RewTerm.weight| & 任务基类经验值\,\cite{rudin2021minutes} & 调高某项使策略更重视它，过高会压制其他项 & 与其他权重\emph{相对}大小才有意义；改一项常需重标其他项 \\
\Verb[breaklines]|EventTerm.mode| & 经核对 Isaac Lab\,\cite{mittal2025isaaclab} & 决定 DR 在 startup/reset/interval 何时触发 & 质量随机化常 \verb|startup|，推力常 \verb|interval|；与 episode 长度耦合 \\
\verb|ObsTerm.func| & \Verb[breaklines]|isaaclab.envs.mdp|\,\cite{mittal2025isaaclab} & 决定该项提取什么物理量 & 必须返回首维 \verb|num_envs| 的张量；与可部署性耦合（见 \cref{ch:rlmc-obsact}） \\
组名 \verb|policy| & locomotion 基类默认 & 决定哪个网络消费这组 obs & Isaac Lab 默认仅 \verb|policy| 组，\verb|critic| 是追加项（详见 \cref{ch:rlmc-obsact}） \\
\bottomrule
\end{tabular}
\end{center}

\begin{insight}[Manager-Based 的价值有两个互补视角]\label{ins:rlmc-eco-manager-value}
\textbf{软件工程视角}：它实现了\emph{关注点分离}（separation of concerns）——每个 Manager 只负责一件事，改 obs 不碰 reward，改 DR 不碰 reset，降低认知复杂度。\textbf{科学研究视角}：它让\emph{消融实验}（ablation）变得廉价——你只改一行 config 就得到一个新实验条件，确保对比公平。对比单体架构：在 \verb|legged_gym| 里加一个 reward 项要改 \Verb[breaklines]|compute_reward()| 内部、可能影响其他项的计算顺序；在 Manager 架构里只需加一行 \verb|RewTerm()|，其他 reward 完全不受影响。这正是 \cref{paper:rlmc-eco-isaacgym} 单体痛点的解药。
\end{insight}

\subsection{代码走读与三框架对比}
\label{sec:rlmc-eco-manager-compare}

mjlab 与 Isaac Lab 能互通的根本原因，是它们共享同一套 Manager-Based 架构。下面把同一个「四足速度跟踪」任务的配置在 mjlab 与 Isaac Lab 中并排（简化版）——\textbf{注意 mjlab 与 Isaac Lab 在命名上的真实差异}：mjlab 用\emph{原生 dataclass 工厂函数} + 普通 \verb|dict| 装 term（无 \verb|@configclass|），term 类用\emph{全名}（无 \verb|ObsTerm|/\allowbreak\verb|RewTerm| 短别名，那是 Isaac Lab 的 import 重命名习惯），奖励/观测函数名也与 Isaac Lab 不同。再放上 UniLab 的等价配置——\textbf{第三段读起来更不一样}：UniLab 不写 Manager 配置树，而是在 \verb|NpEnv| 子类里直接实现，奖励写成 Hydra 标量字典。这种「同一任务、多种架构哲学」的并置，恰是看清 Manager-Based 设计\emph{得失}的最好镜子。

\begin{codebox}[Python]{mjlab:velocity 任务配置（G1，工厂函数 + dict，简化示意）}
from mjlab.managers import ObservationTermCfg, RewardTermCfg   # mjlab 用全名，无 ObsTerm/RewTerm 别名
from mjlab.managers import ObservationGroupCfg

def make_g1_velocity_env_cfg():                # mjlab 用工厂函数返回 @dataclass 实例
    actor_terms = {                            # term 容器是普通 dict（顺序靠插入序）
        "base_ang_vel": ObservationTermCfg(func=mdp.base_ang_vel),
        "joint_pos":    ObservationTermCfg(func=mdp.joint_pos_rel),
        "joint_vel":    ObservationTermCfg(func=mdp.joint_vel_rel),
    }
    observations = {"actor": ObservationGroupCfg(terms=actor_terms,   # 组名叫 actor
                                                 concatenate_terms=True)}
    rewards = {                                # mjlab velocity 奖励真名（非 Isaac 的 *_exp）
        "track_lin_vel": RewardTermCfg(func=mdp.track_linear_velocity,  weight=1.5),
        "track_ang_vel": RewardTermCfg(func=mdp.track_angular_velocity, weight=0.75),
    }
    return ManagerBasedRlEnvCfg(scene=..., observations=observations, rewards=rewards)
\end{codebox}

\begin{codebox}[Python]{Isaac Lab:等价 velocity 任务配置（简化示意）}
@configclass
class AnymalVelocityEnvCfg(LocomotionVelocityRoughEnvCfg):
    scene = InteractiveSceneCfg(
        robot=ArticulationCfg(prim_path="{ENV_REGEX_NS}/Robot", ...),  # USD 资产
        terrain=TerrainImporterCfg(prim_path="/World/ground", ...),
    )
    observations = ObservationsCfg(
        policy=ObservationsCfg.PolicyCfg(                  # Isaac Lab 组名叫 policy
            base_ang_vel=ObsTerm(func=mdp.base_ang_vel),
            joint_pos=ObsTerm(func=mdp.joint_pos_rel),
            joint_vel=ObsTerm(func=mdp.joint_vel_rel),
        ),
    )
    rewards = RewardsCfg(
        track_lin_vel_xy_exp=RewTerm(func=mdp.track_lin_vel_xy_exp, weight=1.5),
        track_ang_vel_z_exp=RewTerm(func=mdp.track_ang_vel_z_exp, weight=0.75),
    )
\end{codebox}

UniLab 的等价配置分两半：一个注册到 registry 的 \verb|dataclass| cfg（持 scene / 周期参数），一个继承 \Verb[breaklines]|LocomotionBaseEnv| 的 env 类（在 \verb|update_state| 里\emph{一处}算 obs+reward+terminated）；奖励则\emph{不在 Python 里}，而在 Hydra YAML 的 \verb|reward.scales| 标量字典。下面是与上两段对应的简化骨架（结构经源码核对，\verb|registry|/\allowbreak\verb|NpEnv|/\allowbreak\Verb[breaklines]|run_reward_dispatch| 见 UniLab \Verb[breaklines]|src/unilab/base/| 与 \Verb[breaklines]|envs/locomotion/common/|）。

\begin{codebox}[Python]{UniLab:等价 velocity 任务（NpEnv 子类 + 标量字典奖励，简化示意）}
from unilab.base import registry
from unilab.base.scene import SceneCfg
from unilab.envs.locomotion.common import rewards
from unilab.envs.locomotion.common.base import LocomotionBaseCfg, LocomotionBaseEnv

@registry.envcfg("Go2JoystickFlat")               # CamelCase = training.task_name
@dataclass
class Go2VelocityCfg(LocomotionBaseCfg):
    scene: SceneCfg = field(default_factory=lambda: SceneCfg(
        model_file="assets/robots/go2/scene_flat.xml"))   # MJCF（非 USD）
    reward_config: RewardConfig | None = None     # 奖励权重经 Hydra 注入

@registry.env("Go2JoystickFlat", sim_backend="mujoco")    # 同名可注册多后端
@registry.env("Go2JoystickFlat", sim_backend="motrix")
class Go2WalkTask(LocomotionBaseEnv):
    @property
    def obs_groups_spec(self):                    # 组名固定 obs/critic（非 policy）
        return {"obs": 49, "critic": 52}          # critic 多 3 维 = 特权 base linvel

    def update_state(self, state):                # 无 ObservationManager/RewardManager:
        linvel, gyro = self.get_local_linvel(), self.get_gyro()   # 一处算齐
        dof_pos, dof_vel = self.get_dof_pos(), self.get_dof_vel()
        reward = self._compute_reward(state.info, linvel, gyro, dof_pos)
        obs = self._compute_obs(state.info, gyro, dof_pos, dof_vel)
        terminated = self._backend.get_sensor_data("upvector")[:, 2] <= 0.5
        return state.replace(obs=obs, reward=reward, terminated=terminated)
\end{codebox}

\begin{codebox}[text]{UniLab:奖励是 Hydra 标量字典（conf/ppo/task/go2\_joystick\_flat/mujoco.yaml）}
# @package _global_
training: {task_name: Go2JoystickFlat, sim_backend: mujoco}
reward:                                  # 对照 mjlab/Isaac 的 RewTerm(weight=...)
  scales: {tracking_lin_vel: 1.0, tracking_ang_vel: 0.2,   # 正权重=目标
           lin_vel_z: -2.0, action_rate: -0.01}            # 负权重=正则惩罚
  tracking_sigma: 0.25
\end{codebox}

mjlab 与 Isaac Lab 结构几乎一一对应，差异集中在命名与资产格式；UniLab 则在\emph{架构维度}就不同（无 Manager 类、奖励是标量字典），故同一行里它常是「另一类东西」而非「换个名字」：

\begin{center}
\small
\begin{tabular}{@{}p{2.0cm}p{3.0cm}p{3.4cm}p{4.0cm}@{}}
\toprule
组件 & mjlab & Isaac Lab & UniLab \\
\midrule
环境基类 & \Verb[breaklines]|ManagerBasedRlEnvCfg| & \Verb[breaklines]|ManagerBasedRLEnvCfg|（及其子类） & \verb|NpEnv|/\allowbreak\Verb[breaklines]|LocomotionBaseEnv|（\textbf{非} Manager 基类） \\
场景 & \verb|SceneCfg|+\verb|EntityCfg| & \Verb[breaklines]|InteractiveSceneCfg|+\Verb[breaklines]|ArticulationCfg| & \verb|SceneCfg|（\verb|model_file|+\Verb[breaklines]|fragment_files|） \\
机器人资产 & MJCF（经 \verb|MjSpec| 合成） & USD（支持 URDF 转换） & MJCF / Motrix XML（\Verb[breaklines]|unilab-import-robot| 转 URDF；无 USD） \\
观测组名 & \verb|actor| / \verb|critic| & \verb|policy|（+可选 \verb|critic|） & \verb|obs| / \verb|critic|（env 侧固定两组） \\
观测/奖励项 & 全名 \Verb[breaklines]|ObservationTermCfg| / \verb|RewardTermCfg|（无短别名） & 习惯重命名为 \verb|ObsTerm| / \verb|RewTerm| & 无 term 类：obs 在 \verb|update_state| 拼，reward 是 \verb|scales| 标量字典 \\
数据访问 & \Verb[breaklines]|env.scene.robot|（属性） & \Verb[breaklines]|env.scene["robot"]|（字典） & 后端 getter（\verb|get_dof_pos|/\allowbreak\Verb[breaklines]|get_local_linvel| 等） \\
配置系统 & 原生 \verb|@dataclass| + 工厂函数 & 自研 \verb|@configclass| 嵌套类 & Hydra \verb|@dataclass| + \verb|conf/*.yaml| 覆盖 \\
\bottomrule
\end{tabular}
\end{center}

经核对 UniLab 源码（\Verb[breaklines]|base/registry.py| 的 \Verb[breaklines]|register_env_config|/\allowbreak\verb|register_env|）：任务用\textbf{双装饰器}注册——\Verb[breaklines]|@registry.envcfg("Name")| 标 cfg、\Verb[breaklines]|@registry.env("Name", sim_backend=...)| 标 env 类，同名可挂多个后端；这与 Isaac Lab 的 \verb|gym.register|（字符串 entry\_point）、mjlab 的 \Verb[breaklines]|register_mjlab_task|（已实例化 cfg 对象）是\emph{三种}不同的注册机制，但都回答同一个问题：「为什么 \verb|train| 给个任务名就能跑」。配置覆盖则走 Hydra 深合并（如 \Verb[breaklines]|algo.num_envs=4096|），路由键 \verb|algo|/\allowbreak\verb|task|/\allowbreak\Verb[breaklines]|training.sim_backend| 是保留字、必须用 flag。

\begin{insight}[两类差异：mjlab$\leftrightarrow$Isaac Lab 是「换名字」，UniLab 是「换架构」]\label{ins:rlmc-eco-pickright}
看上表要分清\emph{两种量级}的差异。mjlab 与 Isaac Lab 之间是\textbf{命名差异}（\verb|actor| vs \verb|policy|、\verb|EntityCfg| vs \Verb[breaklines]|ArticulationCfg|），底层概念完全一致——掌握一个迁移到另一个只需 1--2 天，这正是「先用 mjlab 快速验证、需要视觉/多算法时切 Isaac Lab」之所以低成本的原因。而 UniLab 与前两家之间是\textbf{架构差异}：Manager 配置树 vs \verb|NpEnv| 子类、term 类 vs 标量字典、GPU 物理 vs CPU 物理——迁移要重写环境实现层，不是改命名。这条「两类差异」的区分，正是三框架对比的教学价值所在：前一对让你看清\emph{同一架构的不同命名}（概念稳定），后者让你看清\emph{同一问题的不同架构}（设计权衡）——当一个独立团队\emph{反着}主流范式做（把物理放回 CPU），你会自然追问「它在赌什么、放弃了什么」，这个追问比记住任何 API 都更接近工程本质。
\end{insight}

\begin{pitfall}[「两框架代码可直接复制粘贴」是错的]
API 结构类似，但命名与细节不同：把 mjlab config 原样贴到 Isaac Lab 会报 \Verb[breaklines]|AttributeError|/\allowbreak\verb|KeyError|（如组名 \verb|actor| vs \verb|policy|、数据访问属性 vs 字典）。迁移要按对应表逐项替换，而非复制。这也是为什么 \cref{ch:rlmc-obsact} 会专门厘清「环境侧组名」与「runner 侧 \texttt{obs\_groups} 键」两层命名的混乱。
\end{pitfall}

\subsection{运行验证：怎么确认 Manager 接对了}

改完 Manager 配置后，\emph{先验证再训练}。最廉价的验证是 \cref{ins:rlmc-eco-zerorandom} 的 zero/random 体检，再加一步「打印实际维度」：

\begin{codebox}[Python]{改完 obs/reward 配置后先跑几步，核对维度与 NaN（组名自适应三框架）}
import torch

def sanity_check(env, num_steps=5):
    """改完 Manager 配置后跑 5 步 random agent，核对接线."""
    obs_dict, _ = env.reset()
    # 关键:actor 组名跨框架不同——Isaac Lab 叫 "policy"，mjlab 叫 "actor"，
    # UniLab 叫 "obs".别硬编码 obs_dict["policy"]，否则照搬到 mjlab 会 KeyError.
    actor_key = next(k for k in ("policy", "actor", "obs") if k in obs_dict)
    print("实际可用的 obs 组名:", list(obs_dict.keys()), "-> 用", repr(actor_key))
    for step in range(num_steps):
        actions = torch.rand(env.num_envs,
                             env.action_manager.total_action_dim,
                             device=env.device) * 2.0 - 1.0
        obs_dict, rew, *_ = env.step(actions)
        actor_obs = obs_dict[actor_key]
        assert actor_obs.shape[0] == env.num_envs        # 首维必须是 num_envs
        print(f"[step {step}] obs={tuple(actor_obs.shape)} "
              f"rew_mean={rew.mean():.3f} nan={torch.isnan(actor_obs).any().item()}")
\end{codebox}

预期：obs 形状的首维等于 \verb|num_envs|、第二维等于你手算的观测拼接维度；\verb|rew| 不为 NaN；运行不抛 \verb|KeyError|（组名匹配）。若 obs 维度与手算不符，多半是某个 \verb|ObsTerm| 没生效或拼接顺序有误。

\begin{pitfall}[这段 \texttt{sanity\_check} 直接照搬到 mjlab 会 \texttt{KeyError}]
上面的诊断脚本最初是 Isaac Lab 风格、取 \Verb[breaklines]|obs_dict["policy"]|。但实测三框架的 actor 组名\emph{各不相同}：Isaac Lab 是 \verb|policy|、mjlab 是 \verb|actor|、UniLab 是 \verb|obs|（\cref{sec:rlmc-eco-manager-compare} 的对应表已列）。所以脚本里那行 \Verb[breaklines]|next(k for k in ...)| 不是花架子——它先\textbf{打印当前框架真实存在的组名}再取值，正是「不背书里的名字、跑一次问框架要真值」这一方法论的最小落地。读者最该养成的习惯是：\emph{遇到组名相关的 \texttt{KeyError}，第一反应是} \Verb[breaklines]|print(obs_dict.keys())| \emph{看实际有哪些键，而不是猜}。
\end{pitfall}

\begin{insight}[zero/random 体检：具体该在终端/Viser 里看哪几个量]\label{ins:rlmc-eco-smoke-read}
\cref{ins:rlmc-eco-zerorandom} 说 zero/random 是「免费体检」，但只说「要跑、看什么」还停在概念——落到操作，你应当盯\textbf{三处可观测信号}：
\begin{enumerate}
  \item \textbf{Viser 网页里的姿态}（\Verb[breaklines]|play --agent zero| 时）：浏览器开 \Verb[breaklines]|http://localhost:8080| 看 3D 视图。健康的判据是机器人\emph{稳稳站着、足底贴地、不抖}；若它瘫成一摊或某条腿劈出去，就是默认偏置（offset）配错——回 \cref{ch:rlmc-obsact} 的动作偏置链查 keyframe，而\emph{不}是去调 reward。
  \item \textbf{是否「炸飞」/出 NaN}（\Verb[breaklines]|--agent random| 时）：随机动作下机器人应有明显但\emph{有界}的乱动；若瞬间飞出视野、或上面 \verb|sanity_check| 打印的 \verb|nan=True|，是 action scale 太大或缺裁剪。
  \item \textbf{首行训练日志的三个数}（\Verb[breaklines]|train ... --max-iterations 2| 时）：终端会打印每迭代的 \verb|steps/s| 与各 loss。\emph{第一次跑通只需确认三件事}：\verb|steps/s| 是有限正数（哪怕只有几千，见 \cref{ins:rlmc-eco-throughput}）、\Verb[breaklines]|value loss| 是有限数（不是 \verb|inf|/\allowbreak\verb|nan|）、\verb|reward| 非 NaN。三者俱在即「管线打通」，可以放大 env 数正式训练；任一为 NaN/inf 就先回头查 obs/action，别急着调 PPO 超参。
\end{enumerate}
这三处刚好对应仿真的「调试仪器」角色（\cref{sec:rlmc-eco-roles}）：用零成本的两个 dummy agent，在写一行训练代码之前就把「环境装错/配错」和「算法没调好」分开。
\end{insight}

\subsection{从零起步：搭一个能跑通的最小新任务}
\label{sec:rlmc-eco-fromscratch}

前面都在\emph{读}已有任务的配置；但工程上你迟早要\emph{新建}一个任务。新手最容易犯的错是「对着代码框从一张白纸手写整棵 cfg 树」，结果撞上一串 import 路径与类名坑。真实且省力的路径恰恰相反——\textbf{复制框架自带的最近邻任务，再增量改}。下面给出三框架各自「从注册到跑通」的最小步骤；记住这套骨架，比记任何单个 API 都更接近怎么\emph{自己}起一个任务。

\begin{codebox}[bash]{mjlab:从复制现成任务到跑通新任务的最小步骤}
# 1) 先确认基准任务存在、能跑（别从空白起步，先有一个能跑的锚点）
uv run list-envs | grep Velocity-Flat-Unitree-G1
# 2) 在源码里定位它的 env_cfg（task 的 cfg + 注册都在这一带）
#    mjlab 的 velocity 任务在 mjlab/tasks/velocity/ 下;找到 register_mjlab_task(...)
# 3) 复制该任务目录为 my_velocity/，改两处:
#    (a) register_mjlab_task("Mjlab-MyVelocity-G1", cfg=make_my_cfg()) 换个唯一任务名
#    (b) 在 make_*_cfg() 里改你要动的 term（如给 rewards dict 加一项惩罚）
# 4) 重新 list-envs 确认新名字出现，再 zero/random 体检，最后 2 迭代冒烟
uv run list-envs | grep MyVelocity                 # 注册成功的判据:名字在表里
WANDB_MODE=disabled uv run train Mjlab-MyVelocity-G1 \
    --env.scene.num-envs 64 --agent.max-iterations 2   # steps/s 有限 + reward 非 NaN 即通
\end{codebox}

三框架的「最小新任务」骨架同形、只是注册机制不同（已在 \cref{sec:rlmc-eco-manager-compare} 末尾对齐过三种注册）：\textbf{Isaac Lab} 复制 \Verb[breaklines]|source/.../velocity/config/<robot>/| 下现成任务，在其 \verb|__init__.py| 里 \Verb[breaklines]|gym.register("Isaac-MyTask-v0", entry_point=..., kwargs={"env_cfg_entry_point": MyEnvCfg})|，再 \Verb[breaklines]|./isaaclab.sh -p .../train.py --task Isaac-MyTask-v0 --headless --max_iterations 2| 冒烟；\textbf{UniLab} 则复制 \verb|conf/| 下最近的任务 YAML + 一个 \Verb[breaklines]|@registry.envcfg("MyTask")| 的 dataclass 与 \Verb[breaklines]|@registry.env("MyTask", sim_backend="mujoco")| 的 \verb|NpEnv| 子类，用 \Verb[breaklines]|unilab train --algo ppo --task my_task --sim mujoco algo.num_envs=16| 冒烟。

\begin{insight}[「从最近邻复制」为什么几乎总比「从零手写」对]\label{ins:rlmc-eco-copyfrom}
框架自带任务是\textbf{开发者亲手调通、且与当前版本 API 同步}的——它替你踩平了 import 路径（如 mjlab 的 \Verb[breaklines]|ObservationGroupCfg| 实际在 \Verb[breaklines]|mjlab.managers.observation_manager| 而非你以为的某处）、类名、scene 资产挂载这些「不写错就跑、写错就一行 \verb|ImportError|」的隐性约束。从它出发，你每次只\emph{增量}改一两个 term，出错时变量被限制在你刚动的那几行，二分定位极快；而从空白手写，第一次 \verb|python| 往往先撞一串与你真正想改的逻辑无关的环境坑（典型如 \Verb[breaklines]|ModuleNotFoundError|）。这条「复制最近邻再增量改」是贯穿全书的工程默认动作：\emph{先有一个能跑的，再让它变成你要的}，而不是反过来。
\end{insight}

\begin{practice}[练习·Manager 架构]
\begin{enumerate}
  \item \textbf{[实操]} 按 \cref{sec:rlmc-eco-fromscratch} 的四步，在你装的框架里\emph{复制}一个 velocity 任务、改唯一任务名注册成功（\verb|list-envs| 能看到），再 2 迭代冒烟跑通——体会「从最近邻起步」而非从白纸手写。
  \item \textbf{[实操]} 在 mjlab 或 Isaac Lab 的 velocity 任务里\emph{加一行} \verb|RewTerm|（如关节力矩平方惩罚），跑 50 迭代，对比加项前后的步态/奖励曲线，体会「改一行做一次消融」。
  \item \textbf{[对比]} 对照本节 mjlab 与 Isaac Lab 的两段 config，列出至少三个命名差异，并为每个差异猜一个设计动机。
  \item \textbf{[设计]} 为「崎岖地形四足」任务草拟 Manager 骨架：列出你会用到哪些 Manager、各放哪些代表性 term（obs 至少含本体感知 + 命令，event 至少含一种 DR）。
\end{enumerate}
\end{practice}

% ════════════════════════════════════════════════════════════════
\section{框架选型与迁移}
\label{sec:rlmc-eco-select}

\subsection{全方位对比表}

mjlab 与 Isaac Lab 共享 Manager-Based API、UniLab 走异构 CPU-sim 路线，三者为何并存？因为它们在\emph{物理位置、安装方式、视觉能力、RL 后端、生态}上有本质差异，适合不同场景。下表汇总（数值/版本为 2026 年中快照，时效项见注；UniLab 列经源码与实跑核对）：

\begin{center}
\small
\begin{tabular}{@{}p{1.9cm}p{3.0cm}p{3.2cm}p{4.3cm}@{}}
\toprule
维度 & mjlab & Isaac Lab & UniLab \\
\midrule
物理位置 & MuJoCo Warp（\textbf{GPU}） & PhysX / Newton(Beta)（\textbf{GPU}） & MuJoCoUni / MotrixSim（\textbf{CPU} 多线程） \\
安装 & \Verb[breaklines]|pip install mjlab|（数十秒） & conda+pip+Isaac Sim（数十分钟） & \Verb[breaklines]|uv sync --extra mujoco|（或 \verb|motrix|）；不依赖 Isaac Sim \\
启动延迟 & 数秒 & 数十秒（\textbf{首跑 $\ge 5$ 分钟}编译） & 数秒（CPU 后端，无 Isaac Sim 加载） \\
模型格式 & MJCF & USD（支持 URDF 转换） & MJCF / Motrix XML（URDF 经 \Verb[breaklines]|unilab-import-robot|；无 USD） \\
视觉渲染 & Viser（web，轻量） & Isaac Sim（RTX 光追，工业级） & Viser（web，轻量）；Motrix 原生 renderer 仅回放 \\
RL 后端 & RSL-RL（唯一） & RSL-RL / RL Games / SKRL / SB3 & PPO/APPO/SAC/TD3/FlashSAC（on+off-policy 并重） \\
内置任务数 & 较少（velocity/tracking/lift/cartpole 等，实测约 12 个；\textbf{仅 G1/Go1，无 Go2}） & 较多（loco/manip/dexterous 等） & 26 个注册任务（四足/人形 loco、in-hand、motion-tracking；\textbf{无视觉}） \\
直接访问物理数据 & \checkmark（\verb|mjData|/\allowbreak\verb|mjModel|） & 经 API 间接 & \checkmark（\verb|NpEnvState| + 后端 getter） \\
架构范式 & Manager-Based + CUDA-graph & Manager-Based（双工作流） & 异构 CPU-sim/GPU-policy + 共享内存 IPC（\textbf{非} Manager） \\
RTX 渲染 & 否 & 是 & 否 \\
\bottomrule
\end{tabular}
\end{center}

\begin{note}[两类时效性数据本书不固定]
源稿给了大量 Stars 数（如 Isaac Lab 6.5k/7.3k、mjlab 2.2k）与精确小版本号。\pz{Stars 与小版本号时效性极强、不同日期差异大，本书一律不固定具体数字；需要时请查各仓库当下数据。}本书只固定\emph{稳定}事实：mjlab 轻量单命令安装、Isaac Lab 功能全但安装重、两者 API 趋同。
\end{note}

\subsection{选型决策树}

\noindent\textbf{框架选型（按优先级自上而下，首个命中即选）：}
\begin{enumerate}
  \item \textbf{需要视觉输入（RGB/Depth）？} $\to$ \textbf{Isaac Lab}（RTX 渲染是核心优势；显存按含相机估算）。UniLab 当前无视觉任务、不适用；mjlab 渲染轻量、视觉能力有限。
  \item \textbf{机器只有强 CPU、没有大显存 GPU-sim 后端？} $\to$ \textbf{UniLab}（物理在 CPU 多线程，GPU 只跑策略，不依赖 GPU-sim）。
  \item \textbf{想以异步 off-policy（APPO/SAC/TD3）为主，或要 collector--learner 解耦？} $\to$ \textbf{UniLab}（异构 IPC 让异步采集/学习成为一等公民）。
  \item \textbf{需要对比多种 RL 算法（PPO/SAC/Diffusion）？} $\to$ \textbf{Isaac Lab}（RSL-RL / RL Games / SKRL / SB3 多后端）。
  \item \textbf{需要快速原型 + 轻量部署？} $\to$ \textbf{mjlab}（pip 安装、数秒启动）。
  \item \textbf{任务接触密集（足式 / 灵巧手）？} $\to$ \textbf{mjlab}（MuJoCo 接触精度更好）。
  \item \textbf{以上皆否} $\to$ 选社区 baseline 更多的（多为 \textbf{Isaac Lab}）。
\end{enumerate}

\begin{insight}[决策树第一问为何是「要不要视觉」]\label{ins:rlmc-eco-treeroot}
把「视觉」放在决策树根节点，是因为它同时触发\emph{两个}最贵的约束：(1) 渲染质量——只有 Isaac Sim 的 RTX 光追能给逼真纹理/光照（视觉 sim-to-real 的关键）；(2) 显存——含相机的环境每环境显存暴涨 $10$--$50$ 倍（\cref{ins:rlmc-eco-vram}），直接决定你需要多大的卡、能跑多少并行。一旦确定无视觉，后面的选择（多算法/轻量/接触密集）都是「锦上添花」级别的偏好，没有视觉那样的硬约束。
\end{insight}

\subsection{迁移成本评估}

\begin{center}
\small
\begin{tabularx}{\linewidth}{@{}llL@{}}
\toprule
迁移方向 & 成本 & 主要工作 \\
\midrule
mjlab $\leftrightarrow$ Isaac Lab & 低（1--2 天） & 改 config 格式、\verb|EntityCfg|$\leftrightarrow$\Verb[breaklines]|ArticulationCfg|、模型路径 \\
Manager-Based $\leftrightarrow$ UniLab & 中高（约 1 周） & 架构差异：Manager+term 树 $\leftrightarrow$ \verb|NpEnv| 子类 \verb|update_state| + \verb|scales| 标量字典；reward 函数要重写为 \verb|RewardContext| 形式 \\
Isaac Gym $\to$ mjlab/Isaac Lab & 中（3--5 天） & 单体类拆成 Manager config，须读懂旧 obs/reward 逻辑 \\
非 PyTorch（JAX/MJX）$\to$ 上述 & 高（1--2 周） & 整条训练管线重写 \\
\bottomrule
\end{tabularx}
\end{center}

\begin{pitfall}[选型阶段三个高频踩坑]
\textbf{一·为纯 locomotion 任务硬上 Isaac Lab}：一个真实案例是某同学在 Isaac Lab 上花三天配环境（Isaac Sim 版本冲突 + CUDA 不匹配），最后发现任务是纯四足 locomotion、不需视觉也不需多算法——若一开始用 mjlab，\Verb[breaklines]|pip install| + \Verb[breaklines]|uv run train| 十分钟即可开训。\textbf{二·在 mjlab 上做视觉策略}：mjlab 的 Viser 不支持 RTX 光追，视觉 sim-to-real 的纹理/光照随机化能力有限，视觉密集任务应选 Isaac Lab。\textbf{三·以「Stars 多」论优劣}：Stars 反映社区规模而非代码质量；框架自带的内置任务往往由开发者亲写、质量极高，是最好的入门阅读材料。
\end{pitfall}

\begin{practice}[练习·选型]
\begin{enumerate}
  \item \textbf{[决策]} 为下列课题各选一个框架并说明理由：(a) G1 人形在崎岖地形上的速度跟踪；(b) Franka 机械臂的 RGB 抓取；(c) 对比 PPO/SAC/Diffusion 在四足上的表现。
  \item \textbf{[综合]} 若最终要把策略部署到 Unitree Go2 真机，sim-to-real 管线的哪些环节适合用 mjlab、哪些适合 Isaac Lab？有没有需要\emph{同时}用两者的环节（提示：sim2sim 验证）？
  \item \textbf{[规划]} 为你自己的研究课题制定「主线框架 + 辅助框架」计划：什么场景用主线，什么场景切到辅助？
\end{enumerate}
\end{practice}

% ════════════════════════════════════════════════════════════════
\section{其他框架与工具：知道定位即可}
\label{sec:rlmc-eco-tools}

除三大框架外，你读论文与开源代码时一定会遇到下列名词。无需深入，但要知道定位以免走弯路。

\paragraph{Brax / MuJoCo Playground（JAX 生态）。}基于 JAX 的可微分仿真，物理用 MJX 批量化。优势是端到端可微分；局限是 JAX 与主流 PyTorch RL 库不兼容\,\cite{zakka2025playground}。若论文「trained in MJX」而你用 PyTorch，正确做法是在 mjlab（同为 MuJoCo 物理）或 Isaac Lab 上重实现，而非硬迁 JAX 管线。

\paragraph{MuJoCo Menagerie（模型库，非框架）。}Google DeepMind 维护的高质量 MJCF 机器人模型库（Go1/Go2/G1/H1/ANYmal/Franka 等）。mjlab 直接使用它。需要新机器人 MJCF 时先查 Menagerie——用社区验证过的模型远比从零建模可靠。

\paragraph{RSL-RL（RL 算法库，非框架）。}ETH RSL 与 NVIDIA 共同维护的轻量 RL 库（Schwarke 等，arXiv:2509.10771\,\cite{schwarke2025rslrl}）。经核对，它支持 \textbf{PPO + Student-Teacher Distillation}，外加 RND（好奇心探索）与 Symmetry Augmentation（对称数据增广）\,\cite{schwarke2025rslrl}。它是 mjlab 的\emph{唯一}后端，是 Isaac Lab 的可选后端之一；\textbf{连 UniLab 的 PPO 也复用它}——UniLab 在 RSL-RL 的 \Verb[breaklines]|OnPolicyRunner| 上套了个 \Verb[breaklines]|FinalObservationAwarePPO|（处理终止观测的 bootstrap）。三个架构迥异的框架都接同一个 RSL-RL，这件事本身就值得玩味。

\begin{insight}[三框架都接 RSL-RL，暴露了一个事实标准]\label{ins:rlmc-eco-vecenv}
读源码会发现 mjlab 与 Isaac Lab 的 RSL-RL 适配器（各自的 \Verb[breaklines]|RslRlVecEnvWrapper|）\emph{几乎逐行相同}：都继承 RSL-RL 的 \verb|VecEnv|、都在 \verb|step| 后把 obs 包成 \verb|TensorDict|、都把超时标记 \verb|time_outs| 塞进 \verb|extras| 供 value bootstrap。两个独立团队在 RL 胶水层收敛到几乎一样的接口，说明「\verb|VecEnv| 抽象 + \verb|TensorDict| + \verb|time_outs| 入 extras」已是 GPU 并行 RL env 的\textbf{事实标准}。UniLab 虽然环境侧（CPU、\verb|NpEnv|）完全不同，PPO 路径却同样接 RSL-RL \Verb[breaklines]|OnPolicyRunner|——这进一步印证：\emph{算法库与框架是可独立替换的两层}，framework 怎么变，只要实现这套 VecEnv 契约，就能复用同一个 PPO。这正是 \cref{ins:rlmc-eco-layers} 分层可替换观的最硬证据。
\end{insight}

\begin{insight}[「只有 PPO」对 locomotion 通常不是局限]\label{ins:rlmc-eco-ppo}
RSL-RL 的在线 RL 算法以 PPO 为主，初看像局限，但 PPO 统治 locomotion 有工程理由：(1) locomotion 的 reward 分布在训练中剧变（倒地$\to$站立$\to$行走），off-policy 的 replay buffer 易充满过时数据；(2) PPO 只需前向收集数据，数据收集与梯度更新可完全分离，契合 GPU 并行的数据流；(3) 超参少且有物理直觉；(4) Symmetry Augmentation 只与 on-policy 兼容（在 rollout buffer 里左右镜像每个 transition，等效翻倍 batch）。操作任务里 SAC/Diffusion 可能更合适——要 off-policy 有两条路：Isaac Lab 换 RL 后端（RL Games/SKRL/SB3），或换框架到 \textbf{UniLab}（SAC/TD3/FlashSAC 与异步 APPO 是其一等公民，配共享内存 replay buffer）。这回扣 \cref{sec:rlmc-eco-select} 决策树「要不要多算法 / 要不要异步 off-policy」那两问。
\end{insight}

\begin{note}[在 UniLab 上真上手异步/off-policy：一条可复制命令]
\cref{ins:rlmc-eco-ppo} 把 APPO/FlashSAC 说成 UniLab 的「一等公民」，但只有抽象描述读者不知道怎么真跑。UniLab 的 \verb|conf/| 下确有 \verb|appo|（异步 PPO）与 \verb|flashsac|（off-policy）两套算法配置树，换算法只需改 \verb|--algo| 这一个路由键（其余沿用同一任务）：
\begin{codebox}[bash]{UniLab:把同一个 go2 任务从同步 PPO 切到异步/off-policy}
# 同步 on-policy PPO（默认对照基线）
unilab train --algo ppo     --task go2_joystick_flat --sim mujoco algo.num_envs=16
# 异步 PPO（collector/learner 经共享内存解耦，采集与学习重叠）
unilab train --algo appo    --task go2_joystick_flat --sim mujoco algo.num_envs=16
# off-policy（FlashSAC，配共享内存 replay buffer）
unilab train --algo flashsac --task go2_joystick_flat --sim mujoco algo.num_envs=16
\end{codebox}
三条命令\emph{只有} \verb|--algo| 不同，正是 \cref{ins:rlmc-eco-layers} 「算法层可独立替换」在 UniLab 上的具体落地：环境（CPU \verb|NpEnv|）不动，换的是 GPU 那侧的学习算法。\verb|algo.num_envs| 这种点路径走 Hydra 深合并覆盖（\cref{sec:rlmc-eco-manager} 提过的 \verb|algo|/\allowbreak\verb|task| 是保留路由键）。冒烟时把 \verb|num_envs| 设小（如 16）先确认三条都起得来、reward 非 NaN，再放大 env 数比较异步与同步的真实收益。
\end{note}

\paragraph{WandB / Rerun（实验管理与可视化，非仿真）。}WandB（Weights \& Biases）是实验管理平台，mjlab 原生集成（训练曲线/checkpoint/视频/超参自动上传）。机器人 RL 单个四足任务常有 $30$--$60$ 个可调参数，不用实验管理工具几天后就会忘记「那个成功实验用了什么参数」。Rerun 是时间序列数据流可视化，擅长回放「某个 episode 第 137 步到底发生了什么」。

\begin{insight}[框架-引擎-工具是可独立替换的分层]\label{ins:rlmc-eco-layers}
整个生态可分四层：RL 算法层（RSL-RL/RL Games/SKRL/SB3）、框架层（mjlab/Isaac Lab/UniLab）、物理引擎层（MuJoCo Warp/PhysX/Newton/MJX/MuJoCoUni/MotrixSim）、模型与工具层（Menagerie/USD 资产/WandB/Viser/Rerun）。关键洞察是\emph{每一层都可独立替换}：你能在 Isaac Lab 上把引擎从 PhysX 换成 MuJoCo Warp（经 Newton）而不改 MDP 代码；能在 mjlab 上把算法从 RSL-RL 换成自定义 PPO（实现相同 VecEnv 接口）而不改环境；UniLab 更把「物理引擎层」整个挪到 CPU（MuJoCoUni/MotrixSim）、却仍接同一个 RSL-RL（\cref{ins:rlmc-eco-vecenv}）。这解释了为什么本书选「框架」而非「引擎」或「算法」作教学轴线——框架层是连接算法与引擎的\emph{桥}，是研究者日常交互最多的层级；学会框架层（无论它是 Manager-Based 还是 UniLab 的 \verb|NpEnv| 范式），就同时获得了对上下两层的控制力。
\end{insight}

\begin{pitfall}[四个易混的「是什么」]
\textbf{一·把 RSL-RL 当通用 RL 库}：它在线算法以 PPO 为主，无 SAC/TD3；要 off-policy 须用 Isaac Lab + 其他后端。\textbf{二·把 Menagerie 当框架}：它只是 MJCF 模型库，不提供 \verb|env.step()|。\textbf{三·混淆框架与算法库}：mjlab/Isaac Lab 提供 \verb|env.step()|，RSL-RL/RL Games 提供 \verb|ppo.update()|，两者经 VecEnv wrapper 连接。\textbf{四·忽视 WandB}：实验散落本地几周后必丢，从第一次训练就用 WandB 是零成本高回报的习惯。
\end{pitfall}

% ════════════════════════════════════════════════════════════════
\section*{常见误解汇总}
\label{sec:rlmc-eco-misconceptions}
\addcontentsline{toc}{section}{常见误解汇总}

\begin{center}
\small
\begin{tabular}{@{}p{5.4cm}p{8.4cm}@{}}
\toprule
误解 & 正确理解 \\
\midrule
「仿真越真实越好」 & 足够真实 + 域随机化 $>$ 极致真实；过度精细牺牲吞吐、收益递减（\cref{ins:rlmc-eco-analogy}） \\
「GPU 并行就是核心多所以快」 & 核心优势是数据零搬运，obs/action/梯度全程在显存（\cref{ins:rlmc-eco-zerocopy}） \\
「Isaac Gym 还能用」 & 可读旧代码，新项目应用 Isaac Lab 或 mjlab \\
「MJX = MuJoCo Warp」 & MJX 基 JAX，MuJoCo Warp 基 NVIDIA Warp；精度近、编程模型异 \\
「框架 = 物理引擎」 & 框架 $\ne$ 物理引擎 $\ne$ RL 算法库（\cref{ins:rlmc-eco-layers}） \\
「Stars 多 = 质量好」 & Stars 反映社区规模而非代码质量 \\
「sim-to-real 不可能成功」 & 已有大量成功案例（walk-these-ways/humanoid-gym/ASAP 等） \\
「选一个框架学就够」 & 三框架边际成本低、覆盖互补（\cref{ins:rlmc-eco-pickright}） \\
「三框架都是 Manager-Based」 & mjlab/Isaac Lab 是；UniLab 是异构 CPU-sim、\textbf{非} Manager（\cref{ins:rlmc-eco-manager-granularity}） \\
「GPU 并行是机器人 RL 唯一答案」 & UniLab 故意把物理放回 CPU、用共享内存解耦也成立（\cref{ins:rlmc-eco-heterogeneous}） \\
「RSL-RL 只有 PPO」 & 还有 Distillation + RND + Symmetry（\cref{ins:rlmc-eco-ppo}） \\
「\texttt{sim.step} 与 \texttt{env.step} 一回事」 & 前者只推物理，后者封装整套 MDP（\cref{ins:rlmc-eco-simenv}） \\
\bottomrule
\end{tabular}
\end{center}

% ════════════════════════════════════════════════════════════════
\section*{小结}
\label{sec:rlmc-eco-summary}
\addcontentsline{toc}{section}{小结}

本章是「强化学习运动控制」部的地图章，主线是\textbf{以仿真在 RL 管线中的工程职责为骨架、在每个职责下横向铺三框架对比}。我们先立三角色——物理世界（\verb|sim.step|）、MDP 机器（\verb|env.step|）、调试仪器（ground-truth），它们是衡量任何框架的三把尺子（\cref{ins:rlmc-eco-simenv}）；再论证 GPU 并行的刚需性，其本质是\emph{数据零搬运}而非核心多（\cref{ins:rlmc-eco-zerocopy}），上限是显存（\cref{eq:rlmc-eco-maxenvs}）；接着梳理六阶段演进，提炼出「保留核心优势、消除核心痛点」的统一模式（\cref{ins:rlmc-eco-pattern}）；然后落到框架架构，用并排配置说明 mjlab 与 Isaac Lab 同源命名小异、而 UniLab 是\emph{非} Manager-Based 的异构 CPU-sim 第三路线（\cref{ins:rlmc-eco-manager-granularity}、\cref{ins:rlmc-eco-heterogeneous}），三者却都接同一个 RSL-RL（\cref{ins:rlmc-eco-vecenv}）；最后给出以「是否需要视觉」为根的选型决策树（\cref{ins:rlmc-eco-treeroot}）与分层可替换的生态观（\cref{ins:rlmc-eco-layers}）。本章不教某一行 reward，而是给你一张地图，让后续每一章都能定位在整体结构中。

\paragraph{术语速查。}
\begin{center}
\small
\begin{tabular}{@{}ll@{}}
\toprule
术语 & 含义 \\
\midrule
\verb|sim.step()| / \verb|env.step()| & 物理推进一步 / 框架级 MDP step（含 decimation+obs+reward+reset） \\
decimation（action repeat） & 一个 policy step 内调用 \verb|sim.step()| 的次数 \\
SoA / AoS & Structure / Array of Arrays，GPU 偏好 SoA 以合并访存 \\
CUDA Graph & 预录制的 kernel 启动序列，要求计算图静态 \\
SIMT & 同一指令多线程执行，环境并行的硬件基础 \\
Manager-Based & 把 MDP 各维度拆到独立 Manager 的框架架构 \\
\verb|ObsTerm|/\allowbreak\verb|RewTerm|/\allowbreak\verb|EventTerm| & Manager 里的基本配置单元 \\
\verb|solref| / \verb|solimp| & MuJoCo 接触刚度/阻尼 与 约束阻抗参数 \\
MuJoCo Warp / MJX & MuJoCo 的 NVIDIA-Warp / JAX 两种 GPU 加速实现 \\
Newton / Kamino & 多求解器统一引擎 / 其中处理闭环机构的求解器 \\
\bottomrule
\end{tabular}
\end{center}

\paragraph{知识点总表。}
\begin{center}
\small
\begin{tabular}{@{}clll@{}}
\toprule
编号 & 要点 & 对应节 & 难度 \\
\midrule
1 & 仿真三角色：物理世界/MDP 机器/调试仪器 & \cref{sec:rlmc-eco-roles} & 低 \\
2 & \verb|sim.step| vs \verb|env.step| 是架构理解的起点 & \cref{sec:rlmc-eco-roles} & 低 \\
3 & GPU 并行核心是数据零搬运 & \cref{sec:rlmc-eco-gpu} & 中 \\
4 & SIMT 约束环境设计（用 mask 代替分支） & \cref{sec:rlmc-eco-gpu} & 中 \\
5 & decimation 桥接决策频率与物理精度 & \cref{sec:rlmc-eco-gpu} & 中 \\
6 & 显存估算 $\texttt{max\_envs}\approx(\text{VRAM}-2)\cdot10^3/\text{per\_env\_MB}$ & \cref{sec:rlmc-eco-gpu} & 中 \\
7 & MuJoCo 凸优化软接触 vs PhysX 硬接触 & \cref{sec:rlmc-eco-evolution} & 中 \\
8 & 生态六阶段与「保留-消除」模式 & \cref{sec:rlmc-eco-evolution} & 高 \\
9 & Manager-Based 架构与单体架构的对比收益 & \cref{sec:rlmc-eco-manager} & 高 \\
10 & 三框架配置同源、命名小异 & \cref{sec:rlmc-eco-manager} & 高 \\
11 & 以「是否需要视觉」为根的选型决策树 & \cref{sec:rlmc-eco-select} & 中 \\
12 & 框架-引擎-工具分层可独立替换 & \cref{sec:rlmc-eco-tools} & 低 \\
\bottomrule
\end{tabular}
\end{center}

% ════════════════════════════════════════════════════════════════
\section*{延伸阅读}
\label{sec:rlmc-eco-further}
\addcontentsline{toc}{section}{延伸阅读}

\begin{itemize}
  \item \textbf{Isaac Gym}（Makoviychuk 等，NeurIPS 2021，arXiv:2108.10470）\,\cite{makoviychuk2021isaacgym}——GPU 并行仿真的奠基，教「端到端 GPU 训练为何快」。
  \item \textbf{Learning to Walk in Minutes / legged\_gym}（Rudin 等，CoRL 2021，arXiv:2109.11978）\,\cite{rudin2021minutes}——大规模并行 locomotion 标准范式，教「四足 RL 的 obs/reward/DR 该长什么样」。
  \item \textbf{Isaac Lab}（Mittal 等，arXiv:2511.04831）\,\cite{mittal2025isaaclab}——Manager-Based 架构的设计文档，教「如何把 MDP 拆成可组合的 Manager」。
  \item \textbf{mjlab}（Zakka 等，arXiv:2601.22074）\,\cite{mjlab2026}——轻量框架设计哲学，教「Isaac Lab API 如何搭在 MuJoCo Warp 上」。
  \item \textbf{UniLab}（arXiv:2605.30313，文档 unilabsim.github.io/UniLab-doc）\,\cite{unilab2026}——异构 CPU-sim/GPU-policy 架构，教「为什么把物理留在 CPU、用共享内存 IPC 解耦 collector/learner，是 GPU-sim 主导范式之外的第三条路」。
  \item \textbf{MuJoCo}（Todorov 等，IROS 2012）\,\cite{todorov2012mujoco}——凸优化接触的物理引擎原始论文，教「为什么 MuJoCo 接触稳」。
  \item \textbf{MuJoCo Playground}（Zakka 等，arXiv:2502.08844）\,\cite{zakka2025playground}——MJX + Madrona 批量渲染，教「JAX 生态下的可微分仿真与视觉训练」。
  \item \textbf{RSL-RL}（Schwarke 等，arXiv:2509.10771）\,\cite{schwarke2025rslrl}——RL 算法库设计，教「PPO + Distillation + RND + Symmetry 的工程封装」。
  \item \textbf{ASAP}（He 等，RSS 2025，arXiv:2502.01143）\,\cite{he2025asap}——跨仿真器 delta action model，教「物理引擎选择对 sim-to-real 的实质影响」。
  \item \textbf{Newton}（Linux Foundation 开源项目）\,\cite{newton2025}——多求解器统一引擎，官方文档教「如何一行切换物理后端」。
\end{itemize}

% ════════════════════════════════════════════════════════════════
\section*{故障排查手册}
\label{sec:rlmc-eco-trouble}
\addcontentsline{toc}{section}{故障排查手册}

\begin{center}
\small
\begin{tabular}{@{}p{3.0cm}p{3.4cm}p{6.6cm}@{}}
\toprule
症状 & 可能原因 & 排查 / 相关节 \\
\midrule
\verb|import| 即报 \Verb[breaklines]|undefined symbol| / 段错误 & CUDA/PyTorch 与 Isaac Sim 版本不匹配 & 照官方矩阵重装 torch；别用默认 wheel（见\nameref{sec:rlmc-eco-setup}） \\
zero agent 下机器人瘫软/劈叉 & 默认偏置（offset）配错，非算法问题 & 查默认关节位置/MJCF keyframe（\cref{ins:rlmc-eco-zerorandom}，详见 \cref{ch:rlmc-obsact}） \\
random agent 立即炸飞或出 NaN & action scale 太大或缺裁剪 & 打印 processed action 范围（\cref{ins:rlmc-eco-zerorandom}） \\
足端穿地、整机下沉 & 物理层接触参数（\verb|solref| 太软） & 这是「物理世界」角色问题，非 MDP 层（\cref{ins:rlmc-eco-simenv}） \\
策略不跟速度命令、奖励像随机游走 & MDP 层 obs/reward 配错 & 先查接口再查奖励最后才调 PPO（\cref{ch:rlmc-obsact}） \\
启动报 \verb|KeyError|（组名） & 环境组名与 runner \verb|obs_groups| 不匹配 & 对齐组名（\cref{sec:rlmc-eco-manager}，详见 \cref{ch:rlmc-obsact}） \\
某些 step 突然变慢 & CUDA Graph 因动态分支失效或 DR 首次触发重建图 & 正常现象，非 bug（\cref{sec:rlmc-eco-gpu}） \\
4096 环境 OOM & 显存不足，尤其含相机时 & 用 \cref{eq:rlmc-eco-maxenvs} 重估；视觉任务换大卡 \\
两框架同机器人行为不同 & 物理引擎接触模型差异 & 正常现象，ASAP 量化过\,\cite{he2025asap}（\cref{sec:rlmc-eco-select}） \\
mjlab config 贴到 Isaac Lab 报错 & API 命名差异 & 按对应表逐项替换（\cref{sec:rlmc-eco-manager-compare}） \\
Isaac Lab 3.0 代码与 2.3 不兼容 & 四元数序/warp-native/模块改名 & 见下「版本信息速查」的 3.0 变更表 \\
\bottomrule
\end{tabular}
\end{center}

上表是「症状$\to$原因」的速查；但工程上更要紧的是\emph{怎么自己把症状定位到原因}。下面把两个最高频的故障，演示成一次真实的\textbf{排查命令序列}——核心方法只有一句：\emph{别猜，去问机器要真值}。

\begin{codebox}[bash]{排查示范一·训练中途 OOM:用 nvidia-smi 把「显存不足」坐实并二分}
# 症状:train 跑到第 N 迭代崩在 torch.cuda.OutOfMemoryError
# 1) 另开一个终端，训练时实时盯显存占用（每 0.5s 刷新）
watch -n 0.5 nvidia-smi
#    看 Memory-Usage 那列:是否在崩之前就逼近卡上限（如 4090 的 24GB）？
#    含相机的任务每环境显存暴涨 10-50x（见 ins:rlmc-eco-vram），最易在这里爆
# 2) 二分确认是不是 env 数的锅:把 num-envs 砍半再跑，能过就是显存上限问题
WANDB_MODE=disabled uv run train <你的任务> --env.scene.num-envs 2048   # 砍半试
# 3) 用估算公式反推该机能跑多少（eq:rlmc-eco-maxenvs），把 num-envs 定在其下
#    结论:OOM 不是玄学——nvidia-smi 一看占用曲线、num-envs 一砍二分即可定位
\end{codebox}

\begin{codebox}[Python]{排查示范二·启动报 KeyError（组名）:打印实际组名而非猜}
# 症状:env.step 后取 obs_dict["policy"] 抛 KeyError，或 runner 报 obs_groups 不匹配
# 不要猜「是不是该叫 actor」——直接把框架\emph{真实}给出的键打出来:
obs_dict, _ = env.reset()
print("实际 obs 组名 =", list(obs_dict.keys()))
#  Isaac Lab 多半打印 ['policy', ...];mjlab 打印 ['actor', 'critic'];UniLab ['obs','critic']
#  → 看清后把代码里的硬编码键改成实际存在的那个（参见 sec:rlmc-eco-manager-compare 对应表）
#  若是 runner 侧 obs_groups 不匹配:核对训练配置里的 obs_groups 键名与上面这行一致
\end{codebox}

\begin{insight}[「排查」与「看症状表」的区别：把书里的名字/数字当\emph{待验证假设}]\label{ins:rlmc-eco-debugmethod}
症状表给你\emph{方向}，但真正的排错动作永远是\textbf{让机器输出真值、再与你的假设对照}：OOM 就 \verb|nvidia-smi| 看占用、\verb|num-envs| 砍半二分；\verb|KeyError| 就 \Verb[breaklines]|print(obs_dict.keys())| 看真名；某个 reward 项疑似失效就去训练日志里看对应的逐项回合奖励是不是恒为 0（三框架都按项打印，前缀各异：mjlab/Isaac Lab 形如 \Verb[breaklines]|Episode_Reward/<term>|，UniLab 形如 \verb|reward/<term>|，哪项恒为 0 即定位到失效的 term）。这套「不背书、跑一次问框架」的习惯，比记住任何一张症状表都更通用——因为版本会变、API 名会漂，但「打印真值再对照」这个动作永远有效。本书所有具体 term 名、维度、版本号，都请当成\emph{待你本地实测复核的假设}，而非真理。
\end{insight}

% ════════════════════════════════════════════════════════════════
\section*{版本信息速查}
\label{sec:rlmc-eco-version}
\addcontentsline{toc}{section}{版本信息速查}

本部后续章节的安装命令与 API 调用以下表为「最低兼容参考」（\pz{快照为 2026 年中；具体小版本号与 Stars 等时效数据请以各仓库当下为准}）：

\begin{center}
\small
\begin{tabularx}{\linewidth}{@{}llL@{}}
\toprule
组件 & 锚定版本 & 说明 \\
\midrule
mjlab & 1.x（\pz{2026-05 PyPI 最新 1.4.0\,\cite{mjlab2026}}） & \Verb[breaklines]|pip install mjlab| / \Verb[breaklines]|uv sync| \\
Isaac Lab & 2.3.2（主线） & conda+pip+Isaac Sim；3.0 Beta 变更见下 \\
UniLab & arXiv:2605.30313\,\cite{unilab2026} & \Verb[breaklines]|uv sync --extra mujoco/motrix|；后端 \verb|mujoco-uni|\,3.8.0 / \Verb[breaklines]|motrixsim-core|\,0.8.2 \\
rsl-rl-lib & Isaac 3.0.1 / mjlab 5.2.0 & 新 config（拆分 actor/critic 模型；PyPI 无对外 4.0 版本号）；三框架 PPO 共用 \\
Newton & 1.0 GA & 引擎可用，Isaac Lab 集成仍 Beta \\
MuJoCo Warp & 随 mjlab & Google DeepMind + NVIDIA 共同维护 \\
Python & $\ge 3.10$ & mjlab 支持 3.10--3.13\,\cite{mjlab2026} \\
\bottomrule
\end{tabularx}
\end{center}

\begin{note}[Isaac Lab 3.0 的破坏性变更（升级前必读）]
本书正文以 Isaac Lab 2.3.0 为主线，但你需知道 3.0 Beta（基于 Isaac Sim 6.0 + Newton）的几处破坏性变更，以便升级报错时定位方向：\textbf{(1) 四元数序}从 \verb|wxyz| 改为 \verb|xyzw|——所有硬编码四元数需更新（与全书 Hamilton 四元数约定无关，是 Isaac Lab 内部存储序）；\textbf{(2) 数据管线}从返回 torch tensor 改为返回 \verb|wp.array|（需 \verb|wp.to_torch()| 包装）；\textbf{(3) 模块前缀}从 \Verb[breaklines]|omni.isaac.lab.*| 改为 \verb|isaaclab.*|；\textbf{(4) 安装}新增 kit-less 模式（可不依赖 Isaac Sim，改用 Newton 后端，运行时以 \texttt{presets=newton} 切换）。以上四项均已对照 Isaac Lab 3.0 官方文档与 \texttt{v3.0.0-beta} Release Notes 核实；其中四元数改序是为对齐 Warp/PhysX/Newton 约定，官方另附 quaternion finder 与 \texttt{scripts/\allowbreak tools/\allowbreak wrap\_warp\_to\_torch.\allowbreak py} 等迁移脚本（脚本名与细节随版本，升级时以官方迁移指南为准）。
\end{note}
